\section{F. Splittings and the specified Selmer problems}

\subsection{Which splittings are being computed}

\begin{statement}{F.1}{Proved}{Sections, torsors, and local conditions}
Let $k$ be a number field, $X/k$ a smooth geometrically connected
curve, and $b$ a rational point or rational tangential point.
For a characteristic quotient $U$ of the geometric fundamental
group, the section at $b$ identifies the corresponding
arithmetic extension with $U\rtimes G_k$. Continuous splittings
modulo conjugation by $U$ are the pointed set
$H^1_{\rm cts}(G_k,U)$. For a unipotent $\Q_p$-group $U$,
a Selmer scheme is a specified subset of this pointed set,
cut out by ramification and local conditions. It is not the
unrestricted set of splittings.
\end{statement}
\begin{proof}
A splitting has the form $\sigma\mapsto(c_\sigma,\sigma)$.
Multiplicativity is equivalent to
$c_{\sigma\tau}=c_\sigma\,\sigma(c_\tau)$; conjugation by
$u\in U$ changes $c_\sigma$ to $u^{-1}c_\sigma\sigma(u)$,
with the inverse convention if conjugation is defined on
the other side. These are precisely nonabelian cocycles
and their equivalence. A different reference section twists
the action and gives the corresponding bijection of
pointed sets, with no canonical choice of origin.
Unramifiedness, crystallinity at $p$, and the specified
local Kummer conditions are additional requirements on
these cocycles. They must be supplied before applying a
dimension or finiteness theorem. This is the precise
meaning used below for the repeated study of splittings
in \MS{51}{734;736}{173;178}, equations (371)--(373).
\end{proof}

\begin{statement}{F.2}{Proved}{All standard tripod dimensions through class eight}
Put $X=\mathbb P^1_\Q\setminus\{0,1,\infty\}$ and use the
unit tangent at $0$. Let $S$ be a finite set of rational
primes, $s=|S|$, and choose $p\ge5$ outside $S$.
For $U_n=\pi_1^{\rm un}(X_{\overline\Q},b)/\gamma_{n+1}$
over $\Q_p$, impose unramifiedness outside $S\cup\{p\}$
and crystallinity at $p$. Write $\operatorname{Sel}_{S,n}$
for this standard Selmer scheme, without additional refined
conditions at $S$. Then
\[
\begin{split}
 \dim\operatorname{Sel}_{S,n}
  &=2s+\sum_{\substack{3\le j\le n\\j\ {\rm odd}}}\ell_j(2),\\
 \dim H_f^1(\Q_p,U_n)&=\sum_{j=1}^n\ell_j(2),\\
 \ell_j(2)&=\frac1j\sum_{d\mid j}\mu(d)2^{j/d}.
\end{split}
\]
The first formula includes $n=1$. In particular:
\end{statement}
\begin{center}\begin{tabular}{r|rrrrr}
\toprule Class&Graded rank&Local&Global $S=0$&Global $|S|=1$&Global $|S|=2$\\\midrule
1 & 2 & 2 & 0 & 2 & 4\\
2 & 1 & 3 & 0 & 2 & 4\\
3 & 2 & 5 & 2 & 4 & 6\\
4 & 3 & 8 & 2 & 4 & 6\\
5 & 6 & 14 & 8 & 10 & 12\\
6 & 9 & 23 & 8 & 10 & 12\\
7 & 18 & 41 & 26 & 28 & 30\\
8 & 30 & 71 & 26 & 28 & 30\\
\bottomrule\end{tabular}\end{center}
For the class-$n$ free metabelian quotient $M_n$ the corresponding
exact dimensions, for $n\ge1$, are
\[
 \dim\operatorname{Sel}^{\rm met}_{S,n}
       =2s+k(k+1),\quad
 \dim H_f^1(\Q_p,M_n)=2+\frac{n(n-1)}2,\quad
 k=\left\lfloor\frac{n-1}{2}\right\rfloor.
\]
The rows at $n=1,\ldots,8$ are:
\[
\begin{array}{c|rrrrrrrr}
n&1&2&3&4&5&6&7&8\\\hline
\dim H_f^1(\Q_p,M_n)&2&3&5&8&12&17&23&30\\
\dim\operatorname{Sel}^{\rm met}_{\varnothing,n}
 &0&0&2&2&6&6&12&12
\end{array}
\]
The rows for $s=1,2$ add $2,4$ respectively. The abelian
approximation is the constant class-one row.
\begin{proof}
The central layer of $U_n$ is $\Q_p(n)^{\ell_n(2)}$.
The mixed-Tate Galois category with the stated conditions
has a graded free prounipotent Lie algebra with $s$
generators in weight one and one generator in each odd
weight at least three. Here weights are Tate degrees,
half the absolute value of the usual motivic weight.
The inputs are the Ext computations and freeness in
\citet[(2.1.3), Proposition 2.3]{DeligneGoncharov2005},
and their continuous $\ell$-adic constrained weighted realization in
\citet[Proposition 9.6, proved using Proposition 9.3 and
Theorem 9.5]{HainMatsumoto2003}.
Constrained weighted completion in their \S\S9.3--9.4 imposes
crystallinity. The degree-one group is
$\mathcal O_{\Q,S}^{\times}\otimes\Q_p$, of dimension $s$:
Kummer classes of $p$ itself are excluded by crystallinity.
The higher Ext groups are one-dimensional in odd degree
at least three and zero in even degree. The degree-two
Ext obstruction groups vanish. We use these actual
Galois completion statements, not a conjectural equality
with all of $\GRT$.

We also check the based action, because the outer
crystallinity conclusion in Hain--Matsumoto's Theorem 9.8
alone would not suffice. At the unit tangent the Belyi
normalization is $x\mapsto x^\chi$,
$y\mapsto f^{-1}y^\chi f$ with $f$ in the commutator
subgroup. It gives the unique normalized lift of the
peripheral outer action, by
\citet[Theorem 1]{HarbaterSchneps2000} and the centralizer
normalization used in B.12. At each unipotent level these
are algebraic word and conjugation formulas. Thus the
based representation factors through the same constrained
weighted completion as the outer one. This is the
additional check needed for its path-torsor cocycles.

Here is the deduction of the dimension formula, including
why it is an equality. Choose a Levi splitting and homogeneous
free generators of that Galois Lie algebra. A cocycle to $U_n$
can be uniquely normalized on the Levi $\Gm$, since all
weights of $\Lie U_n$ are strictly negative. It then
amounts to a graded Lie section into
$\Lie U_n\rtimes\mathfrak u_{\rm Gal}$. On a free generator
of degree $j$, its value may be the given generator plus
an arbitrary element of $(\Lie U_n)_j$; freeness supplies
the unique homomorphism for all such choices. Conversely
every normalized cocycle has precisely these values.
Thus the scheme is, noncanonically, an affine space with
two coordinates for each degree-one Galois generator and
$\ell_j(2)$ for its generator in odd degree $j\ge3$.
The choices of generators and Levi splitting explain the
noncanonicity; the dimension and the Selmer functor do not
depend on them. This also proves compatibility of the
formulas with the tower and describes its inverse-limit
Selmer functor as the corresponding product of coordinates.

At $p$, the crystalline comparison identifies
$H_f^1(\Q_p,U_n)$ with $U_n^{\rm dR}/F^0$.
All Hodge degrees of this tripod Lie algebra are negative,
so $F^0$ is trivial. Witt's formula proves the local
dimension. For the free metabelian Lie algebra on two
generators the derived ideal is the free rank-one module
generated by $[x,y]$ under the commuting operations
$\ad x,\ad y$. Consequently its degree-$j$ rank is $j-1$
for $j\ge2$. Substitution into the same cocycle argument
gives both metabelian formulas and their full inverse
tower. The code computes every row by exact integer
arithmetic; it does not enumerate rational or integral
points.
\end{proof}

\begin{statement}{F.3}{Known}{The first dimension drops and integral finiteness}
For each of $U_n$ and $M_n$, the first strict inequality
\[
 \dim\operatorname{Sel}_{S,n}<\dim H_f^1(\Q_p,U_n)
\]
(with $M_n$ substituted in its case) occurs at class
$1,2,4$ for $s=0,1,2$, respectively. Each inequality
implies that $X(\Z[1/S])$ is finite. It does not identify
that set with the corresponding Selmer locus.
\end{statement}
\begin{proof}
Direct comparison of the exact rows of F.2 proves the
minimality of all three thresholds. The abelian
approximation alone gives the strict inequality only
for $s=0$. For the nonabelian towers,
\citet[Theorem 1 and \S3]{Kim2009} supplies Zariski density
of the unipotent Albanese map on each residue disk and
the dimension criterion for finiteness. His notation
$U_2$ begins with abelianization; our $U_1$ is his $U_2$.
We make this indexing shift explicit before invoking
his theorem.

The global-to-local localization morphism has image
closure of dimension at most the global Selmer dimension.
Under the strict inequality there is a nonzero regular
function on the local space vanishing on that closure.
Its pullback by the unipotent Albanese map is a nonzero
Coleman function on each residue disk. The map for the
metabelian quotient is also Zariski dense: it is the
composition of the dense full map and the surjective
quotient morphism. There are only finitely many integral
residue disks at $p$. On each compact disk a nonzero
analytic function has finitely many zeros, so the
global integral points, whose Kummer classes lie in
the localization image, form a finite set. This proves
the manuscript's proposed mechanism of refining
splittings in this explicitly defined tower; it is not
a finiteness theorem for unrestricted sections.
\end{proof}

\begin{statement}{F.4}{Refuted}{Abelian approximations do not separate every cusp}
For the tripod the unrestricted abelian section sets are
\[
 H^1(G_\Q,\widehat\Z(1)^2)
       =\left(\varprojlim_m\Q^\times/(\Q^\times)^m\right)^2,
 \qquad
 H^1(G_\Q,\Q_p(1)^2)
       =\left(\widehat{\Q^\times}^{\,p}\otimes_{\Z_p}\Q_p\right)^2.
\]
They are infinite. The two tangential sections at $0$ and $1$,
using respectively parameters $z$ and $1-z$ with unit tangent,
have the same class in both sets. Thus the assertion that
distinct cusps remain nonconjugate in every abelian
approximation is false. This counterexample concerns an
approximation, and does not identify the full profinite
cuspidal sections.
\end{statement}
\begin{proof}
The finite Kummer exact sequence and Hilbert 90 give
$H^1(G_\Q,\mu_m)=\Q^\times/(\Q^\times)^m$.
Here $\widehat{\Q^\times}^{\,p}$ denotes the power
completion $\varprojlim_j\Q^\times/(\Q^\times)^{p^j}$.
For this infinite-rank group it is not the completion
over all finite $p$-group quotients.
The inverse systems of finite invariant root groups
are Mittag--Leffler, so the $\varprojlim^1 H^0$ term
vanishes and passage to $\widehat\Z(1)$ gives the first
formula. The identical argument for powers of $p$ gives
the integral $p$-adic formula. A continuous cocycle from
compact $G_\Q$ to $\Q_p(1)$ has bounded image, hence after
one common multiplication by a power of $p$ it is
$\Z_p(1)$-valued. This proves the tensor assertion.

These groups are quite explicit. Write
$\Q^\times=\{\pm1\}\oplus\bigoplus_q\Z$ using valuations
at rational primes $q$. In the profinite formula an
element is a sign and a family $a_q\in\widehat\Z$ such
that, for every positive $m$, $a_q$ is zero modulo $m$
at all but finitely many $q$. In the $p$-adic integral
formula the analogous family lies in $\Z_p$ and tends
to zero $p$-adically outside finite sets, together with
the sign torsion if $p=2$. Tensoring permits a common
bounded denominator. This computes the named abelian
cohomology groups, including their possible infinite
ramification; it does not mistake a direct sum for
its completion.

The two coordinates of abelian path torsors are the
Kummer torsors of $z$ and $1-z$. At a tangential
endpoint, use the leading coefficient of a vanishing
function in the chosen parameter. At the unit tangent
of $0$ these two regularized values are $(1,1)$;
at the unit tangent of $1$ they are again $(1,1)$.
All their Kummer classes vanish. Hence both section
classes are the origin relative to the first section.
Equivalently the connecting path has zero abelian
Galois cocycle. The corresponding standard Selmer
classes for $S=\varnothing$ also coincide, consistently
with the dimension zero in F.2.
\end{proof}

\begin{statement}{F.5}{Proved}{The effect of the degree-seven nonpentagonal element}
Let $\ell$ be any prime and let $\beta$ be the pro-$\ell$
element supplied by B.12 for $\ell\ge5$ and by B.14d
for $\ell=2,3$, with normalized conjugator $f$ having
leading Lie term $q_7$. It acts identically on the
integral lower-central quotients of class at most seven,
and hence on the corresponding $\mathbb Q_\ell$-unipotent
quotients $U_n$ for $n\le7$.
In class eight, with the normalization
$\beta(x)=x,\ \beta(y)=f^{-1}yf$, its first change is
\[
 \beta(\log x)=\log x,\qquad
 \beta(\log y)=\log y+[\log y,q_7]\pmod{\text{degree}>8}.
\]
It transports cocycles for $\rho$ bijectively to cocycles
for $\rho^\beta=\beta\rho\beta^{-1}$. It need not act
on the section set for the \emph{fixed} action $\rho$.
For the primes $p\ge5$ and the local conditions specified
in F.3, all tripod finiteness thresholds occur before
this change can be detected.
\end{statement}
\begin{proof}
Existence for every prime follows from B.12 and B.14d.
The group commutator with an element of $\gamma_7$
lies in $\gamma_8$, proving the integral identity
through class seven without any denominator assumption.
Over $\mathbb Q_\ell$, conjugation by
$\exp(q_7+\text{higher terms})$ changes
a degree-one generator first in degree eight, and its
first term is the displayed bracket. This proves the
identity through class seven and the explicit class-eight
action; it also descends to every characteristic quotient.
If $c_{\sigma\tau}=c_\sigma\rho(\sigma)c_\tau$, application
of $\beta$ gives the cocycle equation for $\rho^\beta$.
Coboundaries are transported in the same way, proving
the bijection. For a fixed action, the same calculation
requires $\beta\rho(\sigma)=\rho(\sigma)\beta$, or an
additional specified intertwining cocycle; the two
cycle equations alone do not give that condition.
Selmer conditions are preserved only when their local
structures are transported as well. In particular a
non-arithmetic element cannot be used to change a count
of fixed-action Selmer classes by this argument.
At classes $1,2,4$ it acts identically, so no such
change is present at the computed dimension-drop levels.
\end{proof}

\subsection{Part F(c). Two punctured elliptic curves, with exact local dimensions}

This subsection implements the successive examination of splittings proposed
in the closing pages of \S51, folio 734, scan 173, and folio 736, scan 178.
It concerns the geometric group of type $(1,1)$ occurring in \S52,
folios 743--746, scans 191--197. It supplies calculations for the abelian,
nilpotent and free metabelian unipotent approximations, with their Galois
actions. Splittings always mean conjugacy classes of splittings satisfying
the specified local conditions. Their parameter schemes can have positive
dimension even when the set of integral points is finite.

Let $E_d:y^2=x^3-d^2x$, $X_d=E_d\setminus\{O\}$, and take $d=1,6$.
We use $p=5$ and the rational base point $b=(0,0)$. Let $\mathcal E_d$
be the minimal proper regular model over $\mathbb Z$ and let
$\mathcal X_d=\mathcal E_d\setminus\overline{\{O\}}$. This model contains
the integral affine points on the displayed Weierstrass equation. We allow
a finite set $S$ of denominator primes, with $5\notin S$, and put $s=|S|$.
The bad primes are $B_1=\{2\}$ and $B_6=\{2,3\}$; put
$T=S\cup B_d\cup\{5\}$. They are not automatically counted in $s$:
at a bad prime outside $S$ the integral local condition is retained.

Write $U_n=\pi_1^{\rm un}(X_{d,\overline{\mathbb Q}},b)/\gamma_{n+1}$,
over $\mathbb Q_5$, and $V=T_5E_d\otimes\mathbb Q_5$. Thus $U_1=V$.
Write $M_n$ for the corresponding class-$n$ quotient of the free metabelian
quotient $U/[U',U']$. Both towers have $L_1=V$, $L_2=\mathbb Q_5(1)$.
For $U$ the higher central layers are $L_n=\operatorname{Lie}_n(V)$;
for $M$ write $L_n^{\rm met}$.
For each finite reduction type $\Sigma$, impose at $\ell\ne5$ the
Zariski closure of the local Kummer image of the indicated integral
component, or of the indicated punctured component at $\ell\in S$.
At $5$ impose crystallinity. Denote the resulting finite union of
refined Selmer schemes by $\operatorname{Sel}_{S,n}$, or
$\operatorname{Sel}^{\rm met}_{S,n}$. Its dimension means the maximum
dimension of its nonempty components. This convention permits the cuspidal
local component at every prime in $S$, even when it has no global integral
point. These schemes still contain the Kummer images of all
$\mathcal X_d(\mathbb Z[1/S])$.

\begin{statement}{F.6}{Proved}{The examples have ranks zero and one, and the usual Selmer ranks agree}
The curves $E_1$ and $E_6$ have Mordell--Weil ranks $0$ and $1$, respectively.
They both have good ordinary reduction at $5$, with $a_5=-2$.
Their Tate--Shafarevich groups are finite; consequently
\[
 \dim H_f^1(\mathbb Q,V)=r,\qquad r=0\ \hbox{or}\ 1,
 \qquad \dim H_f^1(\mathbb Q_5,V)=1.
\]
Here $H_f^1$ uses the Bloch--Kato local conditions at every place.
\end{statement}
\begin{proof}
The pure $2$-Selmer rank of a squarefree congruent-number curve is the
nullity of its Monsky matrix. In the additive Legendre-symbol convention,
for $d=2^\epsilon P_1\cdots P_t$ the even matrix is
\[
 M_e=\begin{pmatrix}D_2&A+D_2\\ A^{\mathsf T}+D_2&D_{-1}\end{pmatrix},
 \qquad
 A_{ij}=\left[\frac{P_j}{P_i}\right]\ (i\ne j),\quad
 A_{ii}=\sum_{j\ne i}A_{ij},
 \quad (D_h)_{ii}=\left[\frac h{P_i}\right].
\]
The empty odd matrix for $d=1$ has nullity zero. For $d=6$, $t=1$,
$P_1=3$, $A=0$, $D_2=D_{-1}=(1)$, so
$M_e=\left(\begin{smallmatrix}1&1\\1&1\end{smallmatrix}\right)$ has
nullity one. Thus the ranks are at most zero and one. This is the
explicit descent formula in \citet[\S2, (2.1)--(2.3)]{FE_DasSaikia2024},
attributed there to Monsky's appendix to Heath--Brown.

For $E_6$, $P=(12,36)$ satisfies
\[
 2P=(25/4,-35/8),\qquad
 4P=(1442401/19600,1726556399/2744000),\qquad 8P\ne O.
\]
Direct counts give $|E_6(\mathbb F_5)|=|E_6(\mathbb F_7)|=8$.
The kernel of reduction at a good prime is a pro-prime formal group, so
reduction is injective on torsion of order prime to that prime. Combining
the two primes shows that the order of rational torsion divides $8$.
Thus $P$ is not torsion, and the rank is one. The same counts at $3,5$
for $E_1$ are $4,8$; its four visible rational $2$-torsion points therefore
exhaust $E_1(\mathbb Q)$. The counts at $5$ also give $a_5=6-8=-2$ for
both curves, proving ordinarity.

The $2$-descent alone does not prove finiteness of the $5$-primary
Tate--Shafarevich group. We use a separate verified input. The primary
database entries $32.a3$ and $576.c3$ record analytic ranks zero and one
\citep{FE_LMFDBE1,FE_LMFDBE6}. Their conductors are $32$ and $576$.
The database's reliability statement specifies that analytic ranks at
most three at conductor at most $500000$ are rigorous modular-symbol
computations \citep{FE_LMFDBReliability}. Kolyvagin's own
\citet[Theorem 1, printed p.~4; also Example 1, p.~11]{FE_Kolyvagin1990}
states that, for a modular elliptic curve of analytic rank at most one,
the algebraic rank agrees and the Tate--Shafarevich group is finite.
The hypotheses hold here. The Kummer exact sequence after tensoring
with $\mathbb Q_5$ gives the global assertion; the formal logarithm gives
the one-dimensional local assertion.

The program \texttt{code/f\_elliptic\_dimensions.py} checks both descent
matrices, all displayed point operations, and all four reduction counts
by exact rational or finite-field arithmetic. The analytic-rank input is
identified separately in its output and is not purported to be recomputed.
\end{proof}

\begin{statement}{F.7}{Proved}{Full and metabelian local dimensions through class eight}
For either curve the exact local dimensions are
\[
 \begin{array}{c|rrrrrrrr}
 n&1&2&3&4&5&6&7&8\\\hline
 \dim H_f^1(\mathbb Q_5,U_n)&1&2&4&7&13&22&40&70\\
 \dim H_f^1(\mathbb Q_5,M_n)&1&2&4&7&11&16&22&29
 \end{array}
\]
The complete linear data used below are
\[
 \begin{array}{c|rrrrrrrr}
 n&1&2&3&4&5&6&7&8\\\hline
 w_n=\dim L_n&2&1&2&3&6&9&18&30\\
 \operatorname{tr}(c\mid L_n)&0&-1&0&-1&0&-1&0&-2\\
 \dim L_n^{c=-1}&1&1&1&2&3&5&9&16\\
 \dim L_n^{\rm met}&2&1&2&3&4&5&6&7\\
 \dim (L_n^{\rm met})^{c=-1}&1&1&1&2&2&3&3&4
 \end{array}.
\]
The degree is the lower-central degree, equal here to the negative
geometric weight. It is not the doubled Tate grading of the tripod.
\end{statement}
\begin{proof}
The topological fundamental group of a once-punctured torus is free on
two generators. The central graded Lie algebra of its unipotent
completion is therefore the free Lie algebra on $V$. The Hodge
filtration on $V$ has one generator in degree $0$ and one in degree
$-1$. A Lie word of length at least two in just the degree-zero
generator vanishes. Hence $F^0L_n=0$ for $n\ge2$, whereas
$\dim F^0V=1$. The de Rham comparison therefore gives
\[
 \dim(U_n^{\rm dR}/F^0)=\sum_{m=1}^n w_m-1,
 \qquad w_m=\frac1m\sum_{e\mid m}\mu(e)2^{m/e}.
\]
There are no Frobenius-fixed vectors in the pure negative-weight
layers. The crystalline torsor classification identifies the displayed
de Rham quotient with $H_f^1(\mathbb Q_5,U_n)$; alternatively its
dimension follows by successive central extensions and
\citet[Corollary 3.1.11]{FE_Betts2023}.

Complex conjugation on $V$ has eigenvalues $1,-1$. The character form
of the Witt formula is
\[
 \operatorname{tr}(c\mid L_m)
 =\frac1m\sum_{e\mid m}\mu(e)
       \bigl(\operatorname{tr}(c^e\mid V)\bigr)^{m/e},
 \qquad
 \dim L_m^{c=-1}=\frac{w_m-\operatorname{tr}(c\mid L_m)}2.
\]
For the free metabelian Lie algebra the elements
$\operatorname{ad}_x^a\operatorname{ad}_y^b[x,y]$, $a+b=m-2$,
form a basis in degree $m\ge2$. Thus the dimension is $m-1$ and the
negative conjugation eigenspace has dimension $\lfloor m/2\rfloor$.
The Hodge argument is unchanged.

As an independent finite check, the program exhausts all binary words
of each length $1$ through $8$, selects the Lyndon words, and counts
them and their signs $(-1)^{\#y}$. It compares these counts with the
Witt formulas and with the explicit metabelian basis. It also verifies
every CM bidegree multiplicity. This is exhaustive at the stated word
lengths, and does not sample Galois cohomology.
\end{proof}

\begin{statement}{F.8}{Proved}{Exact low-level global dimensions and unconditional finiteness}
For either curve, and every finite $S$ with $5\notin S$,
\[
 \dim\operatorname{Sel}_{S,1}=r,
 \qquad \dim\operatorname{Sel}_{S,2}=r+s.
\]
The same identities hold for the metabelian tower. For $E_1$ the first
strict global--local dimension inequality is at class one, for all
$s=0,1,2$. For $E_6$ and $s=0$ it is at class two. These inequalities
prove finiteness of the corresponding integral sets by the unipotent
Albanese method. In particular
$X_1(\mathbb Z[1/S])=\{(-1,0),(0,0),(1,0)\}$ for every $S$.
\end{statement}
\begin{proof}
At every $\ell\ne5$, potential good reduction gives
$H^1(\mathbb Q_\ell,V)=0$: $H^0$ and $H^2$ vanish by weights and local
duality, and the local Euler characteristic is zero. The ordinary
abelian Selmer condition is thus independent of $S$, and F.6 proves
the first identity.

For class two, the involution $[-1]$ fixes $b$ and commutes with Galois.
It acts by $-1$ on $V$ and $+1$ on the central line. Its eigenspaces
therefore give a Galois-equivariant vector-space splitting
\[
 \operatorname{Lie}(U_2)=V\oplus\mathbb Q_5(1),
 \qquad [(v,z),(v',z')]=(0,v\wedge v').
\]
In these coordinates a cocycle consists of $(a,b)$ with
$da=0$ and $db=-\tfrac12a\mathbin{\cup}a$.
For rank zero only $a=0$ occurs. For rank one take the cocycle of the
rational point $P=(12,36)$ and scale its coordinates by
$(a,b)\mapsto(ta,t^2b)$. This supplies an algebraic lift for every
$t\in\mathbb Q_5$ in the full abelian Selmer line; it remains
crystalline at $5$, since $P$ is integral there and the scaling is an
endomorphism of the class-two Galois group object. Consequently the
global class-two fibres before imposing the remaining local conditions
are torsors under
\[
 H^1(G_{\mathbb Q,T},\mathbb Q_5(1))
   =\mathbb Z[1/T]^\times\otimes\mathbb Q_5.
\]
There is no stabilizer quotient: the relevant twisted $H^0$ groups
vanish, since their central graded pieces have negative weights.

The valuation map from this unit space to
$\bigoplus_{\ell\in T}\mathbb Q_5$ is an isomorphism. Crystallinity
at $5$ fixes the $5$-valuation. At $\ell\in B_d\setminus S$, the
integral local Kummer images are finite, and fixing a reduction type
fixes the corresponding central valuation. At a prime in $S$ the
cuspidal component has Zariski closure the full affine central line:
the puncture is rational, and its leading commutator is the nonzero
generator of $\mathbb Q_5(1)$, so its unbounded integral contact orders
are Zariski dense in that line. Thus the type which is cuspidal at
every prime in $S$ leaves exactly $s$ unit parameters. Types which
are noncuspidal at some of them leave fewer parameters. The unit
valuation isomorphism both satisfies the required constants at bad
primes and proves that the maximal dimension is exactly $r+s$.
The finiteness of the non-$5$ integral local images is
\citet[Theorem 0.1 and Corollary 0.2]{FE_KimTamagawa2008}.

F.7 gives the local dimensions $1,2$ at these levels. A proper
algebraic closed subset of the local period space is cut out by a
nonzero regular function. The pullback is a Coleman function; the
unipotent Albanese map is Zariski dense on each residue disk, so it
does not vanish identically there. On the compact integral disks a
nonzero such analytic function has finitely many zeroes. There are
finitely many disks and finitely many reduction types. This proves
finiteness. The original density theorem is
\citet[Theorem 1, p.~93; finiteness deduction in \S3]{FE_Kim2009}.
His indexing starts with $U_2$ abelian, so his $U_{n+1}$ is our $U_n$.
The elliptic rank-at-most-one application appears explicitly in
\citet[Corollary 0.3 and its proof in \S2]{FE_KimTamagawa2008}, whose
$U_3$ is our class two. The modern refined formulation also follows
from \citet[Theorem A, Remark 1.4, and Lemma 3.2]{FE_LLM2023}.

For rank one the abelian localization is nonzero: a non-torsion
rational point has nonzero elliptic $5$-adic logarithm. Thus class one
does not give a dimension inequality. At $s=0$ class two does. The
description of $E_1(\mathbb Q)$ in F.6 gives the last assertion.
\end{proof}

\begin{statement}{F.9}{Proved}{Exact higher elliptic dimensions through class eight}
For $n\ge3$ define the actual finite-dimensional arithmetic invariants
\[
 \delta_n(E)=\dim H_f^1\bigl(\mathbb Q,L_n^*(1)\bigr),\qquad
 \delta_n^{\rm met}(E)=
       \dim H_f^1\bigl(\mathbb Q,(L_n^{\rm met})^*(1)\bigr).
\]
The exact strict global graded dimensions are
\[
 \dim H_f^1(\mathbb Q,L_n)=\dim L_n^{c=-1}+\delta_n(E),
 \quad
 \dim H_f^1(\mathbb Q,L_n^{\rm met})
       =\lfloor n/2\rfloor+\delta_n^{\rm met}(E).
\]
Consequently for $n\ge2$ there are the unconditional bounds
\[
 \begin{split}
 \dim\operatorname{Sel}_{S,n}
   &\le r+s+\sum_{m=3}^n\bigl(\dim L_m^{c=-1}+\delta_m(E)\bigr),\\
 \dim\operatorname{Sel}^{\rm met}_{S,n}
   &\le r+s+\sum_{m=3}^n
           \bigl(\lfloor m/2\rfloor+\delta_m^{\rm met}(E)\bigr).
 \end{split}
\]
If the relevant defects vanish through the displayed level, these
bounds are equalities. The later calculation F.18 proves
$\delta_n(E_d)=\delta_n^{\rm met}(E_d)=0$ for both curves and
$3\le n\le8$. Therefore the exact global dimensions through
class eight are unconditionally
\[
\begin{array}{cc|rrrrrrrr}
 \multicolumn{2}{c|}{\text{full free tower}}&1&2&3&4&5&6&7&8\\
 r&s&&&&&&&&\\\hline
 0&0&0&0&1&3&6&11&20&36\\
 0&1&0&1&2&4&7&12&21&37\\
 0&2&0&2&3&5&8&13&22&38\\
 1&0&1&1&2&4&7&12&21&37\\
 1&1&1&2&3&5&8&13&22&38\\
 1&2&1&3&4&6&9&14&23&39
\end{array}
\]
\[
\begin{array}{cc|rrrrrrrr}
 \multicolumn{2}{c|}{\text{metabelian tower}}&1&2&3&4&5&6&7&8\\
 r&s&&&&&&&&\\\hline
 0&0&0&0&1&3&5&8&11&15\\
 0&1&0&1&2&4&6&9&12&16\\
 0&2&0&2&3&5&7&10&13&17\\
 1&0&1&1&2&4&6&9&12&16\\
 1&1&1&2&3&5&7&10&13&17\\
 1&2&1&3&4&6&8&11&14&18
\end{array}
\]
The first strict dimension inequalities for rank one are at classes
$2,3,4$ for $s=0,1,2$, unconditionally. The class-two conclusion
is F.8; the later exact calculations F.16 and F.17 prove
$\delta_3(E_d)=\delta_4(E_d)=0$ for both curves and establish
the remaining two first thresholds. The full and metabelian groups
agree through class four. F.18 also evaluates the defects in
degrees five through eight, so every entry of both tables is exact
with the refined local conditions specified above.
\end{statement}
\begin{proof}[Proof of the structural reductions and exact tables]
For a geometric representation $W$, Poitou--Tate duality gives
\[
 h_f^1(\mathbb Q,W)=h^0(\mathbb Q,W)
 +h_f^1(\mathbb Q,W^*(1))-h^0(\mathbb Q,W^*(1))
 +h_f^1(\mathbb Q_5,W)-\dim W^{c=1}.
\]
This precise formula is recorded in
\citet[(5.6)]{FE_LLM2023}. For $W=L_n$, $n\ge3$, both global
$H^0$ terms vanish by weights, and $h_f^1(\mathbb Q_5,W)=\dim W$
by F.7. This proves the exact graded formula, including its defect.
For $L_2=\mathbb Q_5(1)$ the dual invariant term is one and the strict
global group is zero; it cannot be treated by the $n\ge3$ formula.
Allowing denominator primes supplies the class-two unit term $s$ of F.8.

Here potential good reduction is useful beyond dimension estimates.
For $n\ge3$ and $\ell\ne5$,
\[
 H^i(\mathbb Q_\ell,L_n)=0\quad(i=0,1,2).
\]
Indeed after a finite extension the representations are unramified and
pure of weight $-n$; the dual twist has weight $n-2$. Neither weight
is zero. Thus the two invariant groups vanish, and local duality and
Euler characteristic give the assertion. The same argument works for
the metabelian layers. At $5$, $F^0L_n=0$ and the dual twist has
$F^0$ equal to its whole de Rham space, with no invariant vectors.
It follows that
\[
 H^1(\mathbb Q_5,L_n)=H_f^1(\mathbb Q_5,L_n),\qquad
 H_f^1(\mathbb Q_5,L_n^*(1))=0.
\]
Poitou--Tate therefore identifies
\[
 H^2(G_{\mathbb Q,T},L_n)^*
       =H_f^1(\mathbb Q,L_n^*(1)).
\]
There is no dependence on the extra finite primes of $T$ in these
groups: the displayed local groups at such primes vanish.

In a central extension
$1\to L_n\to U_n\to U_{n-1}\to1$, the obstruction to lifting a
Galois cocycle lies in this $H^2$, and every nonempty fibre is a
torsor under $H^1(G_{\mathbb Q,T},L_n)$, since the twisted $H^0$
stabilizers vanish. This proves the unconditional upper bound by
induction. If $\delta_n=0$, all lower global classes lift. At a
non-$5$ place their lift is unique. At $5$ the crystalline condition
also lifts: the de Rham quotients
$U_n^{\rm dR}/F^0\to U_{n-1}^{\rm dR}/F^0$ are surjective by
strictness of the Hodge filtration, and the whole central $H^1$ is
crystalline. Thus it is the fibre of the local lift.
Consequently no additional local equation appears in degree $n$.
The map of refined Selmer schemes is then a torsor, locally on the
base, under the vector group of dimension $\dim L_n^{c=-1}$.
Starting with the maximal class-two component of dimension $r+s$
from F.8 proves equality, rather than merely a conditional upper
bound. The identical argument proves the metabelian assertion.
Summing the rows of F.7 gives the two tables.

For completeness the dual Selmer groups have an explicit CM description.
Over $K=\mathbb Q(i)$ write $V=\psi\oplus\overline\psi$ and
$\psi\overline\psi=\chi_5$. The multiplicity of
$\psi^a\overline\psi^b$ in $L_{a+b}$ is
\[
 \ell_{a,b}=\frac1{a+b}
   \sum_{e\mid\gcd(a,b)}\mu(e)
                  \binom{(a+b)/e}{a/e},\qquad a,b\ge1.
\]
It is one in the metabelian layer. The dual-twist characters are
$\psi^{1-a}\overline\psi^{1-b}$, with these multiplicities;
restriction to $K$ and the order-two descent involution recover the
groups over $\mathbb Q$. The program lists every pair $(a,b)$ through
$a+b=8$. This supplies the finite list of Hecke-character Selmer groups
evaluated in F.18; Hodge weights alone do not force their vanishing.
In particular $L_3=V(1)$, so already
$\delta_3(E)=h_f^1(\mathbb Q,V^*)$ is a distinct problem from
the Mordell--Weil rank calculation.

There are meaningful alternative theorems, whose additional hypotheses
are not needed for the direct calculation in F.18.
\citet[Definition 2.2 and Corollary 2.9]{FE_Bars2007}
proves the requisite second-cohomology vanishing under geometric
regularity of $5$ for $E$ and $K(E[5])$; that regularity has not been
verified here. \citet[Theorem 0.1]{FE_Kim2010CM} proves eventual
global--local dimension inequality unconditionally for a smaller CM
quotient $W$; his Theorem 0.2 gives an explicit bound subject to
specified nonvanishing values of $p$-adic $L$-functions. The quotient
$W$ kills the Lie bidegrees with both entries at least two and is a
quotient of the metabelian group. Thus it gives an alternative
unconditional eventual finiteness argument for both curves, after
enlarging $S$ to include bad primes if necessary. Its existence theorem
does not supply a first level at most eight or the full-table defects.

Every strict inequality in the displayed exact tables implies finiteness
by the same density and analytic argument as F.8; the general effective
refinement is \citet[Theorem 6.2.1(A),(B), and \S6.3]{FE_Betts2023}.
For transfer to LM6, the lemma must retain the specified local
conditions, the rank or finite-$\Sha$ input, the actual Galois action,
and every needed $\delta_m=0$ hypothesis. Replacing these hypotheses
by an assertion about a normalizer Lie algebra supplies no such
arithmetic vanishing theorem.
\end{proof}

\subsection{Exact cyclotomic calculations for all elliptic defects through degree eight}

\begin{statement}{F.16}{Proved}{The first elliptic defects vanish at five}
For $E_d:y^2=x^3-d^2x$, $d=1,6$, at $p=5$, put
$V=T_5E_d\otimes\mathbb Q_5$. Then
\[
 H_f^1(\mathbb Q,V^*)=H_f^1(\mathbb Q,V(-1))=0.
\]
Consequently the full and metabelian class-three defects in F.9
vanish for both curves. Their global dimensions at class three are
$r+s+1$, and their local dimension is $4$. For $E_6$ and $s=1$, class three
is the first strict global--local dimension inequality.
\end{statement}

\begin{proof}
We separate the integral calculation, its analytic interpretation,
and the cohomological deduction. The main-conjecture input is
\citet[Theorem~17.4(1)--(2)]{FA_Kato2004}, including its CM case
proved using Rubin's theorem in \S15. Only height-one primes away
from $5$ are used. The additional large-image hypothesis of
Theorem~17.4(3) is neither assumed nor needed.

Write $[r]_d=\int_r^{i\infty}2\pi i f_d(z)\,dz$, where $f_d$
is the rational newform of $E_d$. At level $32$, use the $48$
Manin generators $[c:d]$, indexed by
\[
 (1,d)\ (0\le d<32),\qquad (c,1)\ (0\le c<32,\ c\text{ even}).
\]
A generator is the path $\{g(0),g(\infty)\}$ for a determinant-one
matrix $g$ with bottom row $(c,d)$. Its relations are
\[
 [c:d]+[-d:c]=0,\qquad
 [c:d]+[d:-c-d]+[-c-d:c]=0.
\]
These present the rational relative modular-symbol space
\citep[Chapter~3, Manin-symbol presentation]{FA_Stein2007}.
Every rational path is reduced by its continued fraction. The
Hecke operator $T_\ell$ is computed on the actual endpoints using
the $\ell+1$ matrices
\[
 \begin{pmatrix}1&j\\0&\ell\end{pmatrix}\ (0\le j<\ell),
 \qquad \begin{pmatrix}\ell&0\\0&1\end{pmatrix}.
\]
The executable \texttt{code/F\_advanced\_symbols.py} gives exact
integer matrices, checks both Hecke actions against every Manin
relation, and computes the following ranks over $\mathbb Q$:
the relation matrix has rank $39$, and the simultaneous equations
for a period functional satisfying $T_3=0$ have rank $46$.
The two resulting primitive integral period functionals form the
columns of a $48\times2$ matrix $\Phi$, recorded completely in
\texttt{F\_advanced\_symbols.json}. They satisfy $T_5\Phi=-2\Phi$.
The rank $46$ and the independence of both columns remain true
over $\mathbb F_5$. Thus their $\mathbb Z_5$-span is the full
kernel of the defining period-functional equations over
$\mathbb Z_5$, rather than a sublattice of unknown $5$-index.

The cusp-form space of level $32$ has dimension one. Its form
corresponds to $E_1$, and $a_3=0$. Consequently this two-dimensional
period quotient is precisely the $E_1$ quotient; integration gives
an injective rational map from it to the two real period directions
in $\mathbb C$. Under the chosen coordinates complex conjugation is
\[
 C=\begin{pmatrix}-1&2\\0&1\end{pmatrix},
 \qquad (1,-1)C=-(1,-1).
\]

Let $\epsilon$ be the primitive quadratic character of conductor
$24$. The elementary Gauss-sum twisting identity gives
\[
 \tau(\epsilon)[r]_6
   =\sum_{v\bmod24}\epsilon(v)[r+v/24]_1,
 \qquad \tau(\epsilon)^2=24.
\]
The right side can be evaluated entirely in the small level-$32$
Manin space. Its exact coordinate pairs are
\[
 \begin{array}{c|rrrrr}
 r&0&1/5&2/5&3/5&4/5\\\hline
 \tau(\epsilon)[r]_6
 &(0,0)&(16,-16)&(0,0)&(0,0)&(-16,16).
 \end{array}
\]
Hence the raw moments $B_j=\sum_{a=1}^4a^j
\tau(\epsilon)[a/5]_6$, for $j=0,1,2,3$, are
\[
 (0,0),\quad(-48,48),\quad(-240,240),\quad(-1008,1008).
\]
Both odd moments reduce to $(2,3)$ modulo $5$.
For $E_1$ the untwisted symbols at the same five arguments are
\[
 (0,1),\quad(2,-3),\quad(2,-3),\quad(-2,1),\quad(-2,1).
\]
Its odd moments $B_1=(-8,-2)$ and $B_3=(-164,64)$ reduce to
$(2,3)$ and $(1,4)$, respectively. Both are nonzero in the odd
period line modulo $5$. These untwisted values and the two
point counts $|E_d(\mathbb F_7)|=8$ are checked by the same program.

Here the normalization does not conceal a period-lattice assumption.
All twisted symbols lie in the displayed integral lattice in
$\mathbb Q^2$. A moment nonzero modulo $5$ in that ambient lattice
cannot be divisible by $5$ in the sublattice generated by the
twisted symbols. In fact the first twisted symbol supplies a
$5$-adic unit multiple of the primitive odd vector $(1,-1)$.
This proves the required unit assertion in the actual odd period
line. No Manin-constant or isogeny-degree conjecture is used.

Let $\alpha\in\mathbb Z_5^\times$ be the unit root of
$t^2+2t+5$. With the odd period line normalized as above, for either
curve the
ordinary modular-symbol measure has values
\[
 \mu(a+5^n\mathbb Z_5)
 =\alpha^{-n}[a/5^n]_d^-
    -\alpha^{-n-1}[a/5^{n-1}]_d^-.
\]
The Hecke relation proves the distribution identity; integrality
of all period coordinates and the unit property of $\alpha$ prove
boundedness. Its Mellin transform is the ordinary $5$-adic
$L$-function. For the Teichmuller character $\omega$, the constant
term on its $\omega^j$ branch reduces, up to a unit, to $B_j$
modulo $5$, because $\omega(a)^j\equiv a^j\pmod5$ and the
constant correction has zero sum against an odd character. In
the twisted example $[0]_6$ itself is zero.
Both odd branches are therefore units in $\mathbb Z_5[[\Gamma]]$,
where $\Gamma=1+5\mathbb Z_5$.

The Tate-twist bookkeeping is essential. Put
$T=T_5E_d(-1)$, so $T\otimes\mathbb Q_5=V^*$ is Kato's
cohomological representation for the weight-two form. By
\citet[Proposition~17.2 and \S17.3]{FA_Kato2004},
\[
 \operatorname{Sel}_\infty(T)
   =\operatorname{Sel}_5(E_d/\mathbb Q(\mu_{5^\infty}))(-1),
 \qquad
 \mathfrak X(T)=\operatorname{Sel}_\infty(T)^\vee.
\]
These are discrete restriction limits and Pontryagin duals,
respectively. Kato's interpolation at $\chi_5\eta^{-1}$
corresponds to the ordinary measure at $\eta^{-1}$
\citep[Theorem~16.2, with $r=1$]{FA_Kato2004}.
Thus the trivial $\Delta=\mathbb F_5^\times$ branch for
$\mathfrak X(T)$ corresponds to the $\omega^{-1}$ branch of
$\mu$. The change in the $\Gamma$ variable is the automorphism
induced by the cyclotomic character and preserves units.
The other choices of rational periods multiply this component
by a nonzero scalar, which is a unit after inverting $5$.
Theorem~17.4(2) therefore implies
\[
 e_1\mathfrak X(T)[1/5]=0.
\]
Indeed $\mathbb Z_5[[\Gamma]][1/5]$ is a regular ring of dimension
one; a finitely generated torsion module vanishing at all its
height-one primes is zero. Finite generation now shows that
$e_1\mathfrak X(T)$, and hence $e_1\operatorname{Sel}_\infty(T)$,
has bounded $5$-power exponent. This conclusion does not assert
integral vanishing of the whole cyclotomic Selmer group.

It remains to justify that this bounds the desired group.
Frobenius at $7$ on either $E_d[5]$ has polynomial $t^2+2$, which is
irreducible over $\mathbb F_5$. Hence $E_d[5]$ is an irreducible
$G_{\mathbb Q}$-module. It has no invariants over
$K_\infty=\mathbb Q(\mu_{5^\infty})$: otherwise those invariants,
being stable under $G_{\mathbb Q}$, would be the whole module,
and its irreducible representation would factor through a finite
quotient of $\mathbb Z_5^\times$. Every irreducible
$\mathbb F_5$-representation of such a quotient is one-dimensional.
The same vanishing holds for $T\otimes\mathbb Q_5/\mathbb Z_5$.
Inflation--restriction therefore gives an injection on torsion
cohomology from $\mathbb Q$ to $K_\infty$.

Fix the finite ramification set $\Sigma=\{2,3,5\}$, together with
the infinite place. All the following global classes belong to
$H^1(G_{\mathbb Q,\Sigma},-)$: the representations are unramified
outside $\Sigma$, and the finite local conditions impose
unramifiedness there. The group
$H^1(G_{\mathbb Q,\Sigma},T)$ is finitely generated over
$\mathbb Z_5$ and its image in rational cohomology is a lattice.
Its intersection with any rational subspace is again a lattice.
This finite-ramification step ensures that reduction to torsion
coefficients cannot erase an entire nonzero rational subspace.

At every place the local finite condition for $V^*$ is zero,
as proved in F.9. Thus a class in $H_f^1(\mathbb Q,V^*)$
maps to a locally trivial torsion class and, after restriction,
to $\operatorname{Sel}_\infty(T)^{\operatorname{Gal}(K_\infty/\mathbb Q)}$.
The image of a nonzero finite-dimensional $\mathbb Q_5$-space
under reduction by its integral cohomology lattice is a nonzero
divisible $5$-primary group. Restriction is injective on this image.
But the target lies in the trivial $\Delta$ component, which
has bounded exponent. A divisible group of bounded exponent is
zero, a contradiction. This proves the vanishing, and the
class-three dimensions follow from F.9.
\end{proof}

\begin{statement}{F.17}{Proved}{The class-four defects vanish at five}
For both curves $E_1$ and $E_6$, at $p=5$,
\[
 \delta_4(E_d)=\delta_4^{\rm met}(E_d)
 =\dim H_f^1(\mathbb Q,\operatorname{Sym}^2V^*)=0.
\]
The exact full and metabelian global dimensions at class four are
$r+s+3$, and the local dimension is $7$.
For $E_6$ the first strict global--local inequalities are therefore
at classes $2,3,4$ for $s=0,1,2$, respectively, unconditionally.
The defects in degrees five and above are separate arithmetic questions.
\end{statement}

\begin{proof}
In degree four the free Lie algebra on a two-dimensional space satisfies
\[
 L_4=\operatorname{Sym}^2V\otimes\det V,
 \qquad L_4^*(1)=\operatorname{Sym}^2V^*.
\]
For example, the equivariant map sends
$(u\mathbin{\odot}v)\otimes(x\wedge y)$ to
$[u,[v,[x,y]]]$. It is symmetric in $u,v$ by Jacobi, since
$[[u,v],[x,y]]=0$, and the three usual degree-four commutators
show that it is an isomorphism. The metabelian relation first
changes the rank-two free Lie algebra in degree five, so the same
layer occurs in both towers. Quadratic twisting does not change
$\operatorname{Sym}^2V^*$; it suffices to treat $E_1$.

Write $K=\mathbb Q(i)$, let $\epsilon_K$ be its quadratic character,
and use $V|_{G_K}=\psi\oplus\overline\psi$ with
$\psi\overline\psi=\chi_5$. Since $5$ splits in $K$, these
characters are defined over $\mathbb Q_5$. There is a decomposition
\[
 \operatorname{Sym}^2V^*
   =W\oplus\mathbb Q_5(-1)\epsilon_K,
 \qquad W=\operatorname{Ind}_{G_K}^{G_{\mathbb Q}}\psi^{-2}.
\]
The quadratic character on the second summand is necessary.
Complex conjugation exchanges the two CM basis vectors and hence
fixes their symmetric product, whereas it acts by $-1$ on
$\mathbb Q_5(-1)$.

The form
\[
 g(z)=\eta(4z)^6=q-6q^5+9q^9+\cdots
\]
has weight three, level $16$, and nebentypus $\epsilon_K$.
Its cohomological $5$-adic representation, in the convention used
by Kato, is exactly $W$.
The identification follows from the square of the CM character
of the level-$32$ form, as given explicitly in
\citet[Theorem~1.1 and the symmetric-square factorization on p.~3]
{FA_LiLongTu2018}. To fix the dual convention, the CM characters
on $V$ are $\psi,\overline\psi$; the automorphic square has
characters $\psi^2,\overline\psi^2$, and the cohomological
representation has their inverses. Thus no additional Tate twist
is inserted into $W$.

We first handle the one-dimensional summand. Its finite local
cohomology is zero at every place, by the argument in F.9 for
the direct summands of $L_4^*(1)$. Soul\'e's finiteness theorem
gives
\[
 H^2(G_{K,\{5\}},\mathbb Q_5(2))=0;
\]
see the author's account
\citet[Corollary~27, with the preceding surjectivity explanation]
{FA_Soule2008}. Only rational vanishing is used.
Adding primes above $2$ and $3$ does not change this group: the
new residue terms involve $\mathbb Q_5(1)$ over finite fields
and vanish rationally, since their Frobenius eigenvalue is not
$1$. Restriction and corestriction, or the two eigenspaces for
$\operatorname{Gal}(K/\mathbb Q)$, then give
\[
 H^2(G_{\mathbb Q,\Sigma},\mathbb Q_5(2)\epsilon_K)=0,
 \qquad \Sigma=\{2,3,5,\infty\}.
\]
The kernel of localization in degree one for the dual twist
$\mathbb Q_5(-1)\epsilon_K$ is dual, by Poitou--Tate, to the
kernel of localization in this zero degree-two group. Its finite
Selmer group is contained in that kernel because all its local
finite conditions are zero. Therefore
$H_f^1(\mathbb Q,\mathbb Q_5(-1)\epsilon_K)=0$.

For $W$, the required integral calculation is
\texttt{code/F\_advanced\_weight3.py}. We give the conventions
and the complete small output needed for the deduction.
Use $24$ representatives of $\mathbb P^1(\mathbb Z/16)$,
each with the polynomial basis $X,Y$, giving $48$ generators.
For a determinant-one matrix $h$ put
\[
 [P,a]h=[h^{-1}P,ah],\qquad
 (hP)(X,Y)=P\bigl(h^{-1}(X,Y)\bigr).
\]
The relations are $I+S$ and $I+R+R^2$, where
\[
 S=\begin{pmatrix}0&-1\\1&0\end{pmatrix},\qquad
 R=\begin{pmatrix}0&-1\\1&-1\end{pmatrix}.
\]
Scaling a bottom row by an odd unit $u$ introduces
$\epsilon_K(u)$. For the Hecke action on the linear polynomials
one uses the integral adjugate; the final coset matrix
$\left(\begin{smallmatrix}\ell&0\\0&1\end{smallmatrix}\right)$
has the additional factor $\epsilon_K(\ell)$.
These are the polynomial Manin-symbol conventions of
\citet[the general modular-symbol chapter]{FA_Stein2007}.

The relation matrix has rank $40$. Adding the equations $T_3=0$
gives rank $46$ and a two-dimensional period quotient; $g$ lies
in it and its two period directions exhaust it. The exact
primitive integral dual basis $\Phi$ is recorded in the JSON
certificate. Both $T_3$ and $T_5$ preserve every relation, and
\[
 T_3\Phi=0,\qquad T_5\Phi=-6\Phi.
\]
The rank $46$ and the column rank two persist modulo $5$.
Consequently $\Phi$ gives the full $\mathbb Z_5$ period lattice
for these equations. Every polynomial path below is evaluated
by an exact continued-fraction reduction.

Write
\[
 I_n(a)=\int_{a/5^n}^{i\infty}
              g(z)(5^nz-a)\,dz
\]
in this period quotient; a fixed nonzero normalization of the
differential is absorbed in the periods. The exact vectors are
\[
 \begin{array}{c|rrrr}
 a&1&2&3&4\\\hline
 I_1(a)&(5,-3)&(-3,-11)&(-11,-3)&(-3,5)\\
 I_0(a)&(1,1)&(1,1)&(1,1)&(1,1).
 \end{array}
\]
Let $\alpha$ be the unit root of $t^2+6t+25$, so
$\alpha\equiv4\pmod5$. The bounded ordinary measure for the
top critical value is
\[
 \mu(a+5^n\mathbb Z_5)
   =\alpha^{-n}I_n(a)-5\alpha^{-n-1}I_{n-1}(a).
\]
The factor $5$ is part of the normalization. Indeed the ordinary
stabilization is $g_\alpha(z)=g(z)-(25/\alpha)g(5z)$;
the change of variable $w=5z$ in $I_n$ turns its second term
into $(5/\alpha)I_{n-1}$. The equation
$U_5g_\alpha=\alpha g_\alpha$ proves the distribution relation.
For an independent finite check, the program evaluates all
$20$ paths with denominator $25$ and verifies exactly, for each
$a=1,2,3,4$,
\[
 \sum_{b=0}^4I_2(a+5b)=-6I_1(a)-5I_0(a).
\]
The integral path coordinates and the unit $\alpha$ prove
boundedness of the measure.

The four residue-class values of $\mu$ modulo $5$ are
\[
 (0,3),\quad(3,1),\quad(1,3),\quad(3,0).
\]
Their Teichmuller moments, for exponents $j=0,1,2,3$, are
\[
 (2,2),\quad(1,4),\quad(4,4),\quad(3,2)
       \quad\text{in }\mathbb F_5^2.
\]
In particular the moment with $j=2$ is a unit in its period
line. As in F.16, the integral ambient lattice makes this a
unit assertion in the lattice generated by the actual modular
symbols; a change of rational periods only multiplies the
associated branch by a nonzero $5$-adic scalar.

Here the relevant branch differs from the weight-two one.
The Fourier expansion and the Gauss-sum identity show that
the Mellin transform at a finite character $\eta^{-1}$ gives,
up to a fixed nonzero period scalar,
\[
 5^{2n}\alpha^{-n}G(\eta)^{-1}L(g,\eta,2)
\]
when $\eta$ has conductor $5^n$. This is Kato's interpolation
at $\chi_5^2\eta^{-1}$ in
\citet[Theorem~16.2, with $r=k-1=2$]{FA_Kato2004}.
Thus the trivial $\Delta$ branch of Kato's $L$-function
corresponds to the $\omega^{-2}$ branch of $\mu$, exactly
the nonzero $j=2$ moment just computed. Twisting the $\Gamma$
variable by $\chi_5^2$ preserves units. Applying
Theorem~17.4(1)--(2) to a stable lattice $T\subset W$ gives
\[
 e_1\mathfrak X(T)[1/5]=0.
\]
As in F.16, finite generation then implies bounded exponent
for the trivial $\Delta$ component of $\operatorname{Sel}_\infty(T)$.
Only the divisibility away from $5$ is used, so no residual
large-image assumption occurs.

There is a minor difference in the final restriction argument:
residual irreducibility is unnecessary here. Let
$A=W/T$ and $H=G_{\mathbb Q(\mu_{5^\infty})}$.
The representation $W$ is irreducible, because the two CM
characters $\psi^{-2}$ and $\overline\psi^{-2}$ are distinct
(their Hodge--Tate weights at either chosen prime above $5$
are distinct). Since $H$ is normal, $W^H$ is either $0$ or $W$.
In the second case $W$ would factor through the abelian
cyclotomic Galois group. Complex conjugation, whose eigenvalues
on this induced representation are $1,-1$, would then commute
with the entire image, making both eigenlines invariant.
This contradicts irreducibility. Therefore $W^H=0$.
The subgroup $A^H$ is finite: it is a subgroup of
$(\mathbb Q_5/\mathbb Z_5)^2$, and positive divisible corank
would give a nonzero invariant Tate module in $W$.
Inflation--restriction shows that the kernel of
\[
 H^1(G_{\mathbb Q,\Sigma},A)\longrightarrow H^1(H,A)
\]
has bounded exponent.

All local finite groups for $W$ are zero by F.9. Hence the
torsion reduction of a global finite Selmer class, restricted
to the cyclotomic extension, satisfies the ordinary kernel
condition in \citet[Proposition~17.2]{FA_Kato2004} and lies in
the trivial $\Delta$ component. As in F.16 we fix
$\Sigma=\{2,3,5,\infty\}$; the finitely generated group
$H^1(G_{\mathbb Q,\Sigma},T)$ supplies a lattice in rational
cohomology. A nonzero rational Selmer subspace would therefore
give a nonzero divisible torsion group. Its restriction lies
in a group of bounded exponent and must vanish; it would
then lie in the bounded-exponent restriction kernel and must
also be zero. This contradiction proves
$H_f^1(\mathbb Q,W)=0$.

Both summands have now been handled, so $\delta_4=0$.
Together with F.16 and the exact lifting argument in F.9,
this gives the class-four global dimension $r+s+1+2=r+s+3$.
The local dimension is $7$ by F.7. For $E_6$ and $s=2$
the dimensions at classes one through four are respectively
$(1,1),(3,2),(4,4),(6,7)$, with the global dimension first.
This proves that class four is the first strict inequality;
the cases $s=0,1$ were settled in F.8 and F.16.
\end{proof}

\begin{statement}{F.18}{Proved}{All elliptic defects through degree eight vanish at five}
For $E_d:y^2=x^3-d^2x$, $d=1,6$, and $p=5$,
\[
 \delta_n(E_d)=\delta_n^{\rm met}(E_d)=0
       \qquad(3\le n\le8).
\]
Consequently every full and metabelian global-dimension entry
in F.9 is unconditional, with the specified refined local conditions.
This conclusion concerns precisely the degrees displayed and
does not require geometric regularity of the entire CM division field.
\end{statement}

\begin{proof}
We first list the representations which must be treated, then
give finite certificates for their required specializations, and
finally prove the passage from those values to the Selmer groups.
Fix $d$, put $K=\mathbb Q(i)$, and write
$V|_{G_K}=\psi_d\oplus\overline\psi_d$.
For $m\ge1$ and $t\ge1$ set
\[
 W_{m,d}=\operatorname{Ind}_{G_K}^{G_{\mathbb Q}}\psi_d^{-m},
 \qquad A_t=\mathbb Q_5(-t),\qquad \epsilon=\epsilon_K.
\]
Suppress $d$ in the next table. The decompositions are
\begingroup\small
\[
\begin{array}{c|l|l}
n&L_n^*(1)&(L_n^{\rm met})^*(1)\\\hline
3&W_1&W_1\\
4&W_2+A_1\epsilon&W_2+A_1\epsilon\\
5&W_3+2W_1(-1)&W_3+W_1(-1)\\
6&W_4+2W_2(-1)+2A_2+A_2\epsilon&W_4+W_2(-1)+A_2\\
7&W_5+3W_3(-1)+5W_1(-2)&W_5+W_3(-1)+W_1(-2)\\
8&W_6+3W_4(-1)+7W_2(-2)+3A_3+5A_3\epsilon
 &W_6+W_4(-1)+W_2(-2)+A_3\epsilon .
\end{array}
\]
\endgroup
Here a plus sign denotes a direct sum. To verify the table, for
$a>b\ge1$ put $m=a-b$ and $t=b-1$. The paired CM weights
$\psi_d^{1-a}\overline\psi_d^{1-b}$ and its conjugate give
$W_{m,d}(-t)$ with multiplicity $\ell_{a,b}$ from F.9; in
particular $m+2t=n-2$. The multiplicity in the metabelian case
is one. In the balanced case $a=b=t+1$, the possible lines
are $A_t$ and $A_t\epsilon$. Their total multiplicity is
$\ell_{t+1,t+1}$, and their difference is
$(-1)^{t+1}\operatorname{tr}(c\mid L_{2t+2})$.
The traces in F.7 give exactly the displayed balanced terms.
These character calculations are valid over $\mathbb Q_5$;
the CM torus and its order-two normalizer act semisimply in
characteristic zero. Both kinds of balanced line have zero
finite Selmer group by the Soul\'e and Poitou--Tate argument
in F.17, now with the positive twist $t+1\ge2$ over $K$.

It remains to treat $W_{m,d}(-t)$ for
\[
 1\le m\le6,\qquad t\ge0,\qquad m+2t\le6.
\]
The cohomological representation $W_{m,1}$ belongs to the
CM form $g_m$ of weight $m+1$ attached to the $m$th power
of the character of $E_1$. The character construction is
reviewed in \citet[\S2.5 and pp.~15--16]{FA_LiLongTu2018}.
For clarity, its level calculation is also explicit. At the
prime $P=(1+i)$ the unit character has order four, with
kernel $1+P^3$ and generator $i$ modulo that kernel.
Its $m$th power therefore has conductor $P^3$ for odd $m$,
$P^2$ for $m\equiv2\pmod4$, and conductor one for
$m\equiv0\pmod4$. Multiplying the conductor norm by the
discriminant $4$ gives the following levels. The central
character is $\epsilon^{m+1}$.
\[
\begin{array}{c|rrrrrr}
m&1&2&3&4&5&6\\\hline
\operatorname{weight}(g_m)&2&3&4&5&6&7\\
\operatorname{level}(g_m)&32&16&32&4&32&16\\
a_5(g_m)&-2&-6&22&-14&-82&234\\
\alpha_m\pmod {25}&13&19&22&11&18&9
\end{array}
\]
The traces in the third row are
$(-1+2i)^m+(-1-2i)^m$. The last row is checked in
$\mathbb Z/25$ against
$\alpha_m^2-a_5(g_m)\alpha_m+5^m=0$ and the unit-root
condition. For $d=6$, odd powers are twisted by the primitive
quadratic character $\chi_{24}$, and even powers are unchanged.
Since $\chi_{24}(5)=1$, the unit roots are unchanged as well.
All these forms are ordinary at $5$.

The program \texttt{code/F\_advanced\_cm\_powers.py}, run with
$m=1,\ldots,6$, constructs the polynomial Manin space of
degree $w=m-1$. It uses the conventions of F.17, with all
$w+1$ monomials $X^{w-j}Y^j$, and isolates the simultaneous
period-functional equations
\[
 R\Phi=0,\qquad T_3\Phi=0,\qquad
 T_5\Phi=a_5(g_m)\Phi.
\]
The CM form supplies two period directions in this kernel.
The exact calculation proves that its dimension is two, so
the kernel is precisely its period quotient. The generator
counts are $48,48,144,24,240,144$; the relation ranks are
$39,40,120,20,200,120$. In every case the simultaneous system
has rank $D-2$ both rationally and modulo $5$, and the two
displayed integral columns have rank two modulo $5$.
The complete matrices $\Phi$ are in the six numbered JSON
certificates. Every Hecke action preserves all Manin relations.

For the top critical value define the integral period vector
\[
 I_n(a)=\int_{a/5^n}^{i\infty}
              g_m(z)(5^nz-a)^w\,dz.
\]
As before, a fixed nonzero normalization is absorbed in the
periods. Its bounded ordinary measure is
\[
 \mu_m(a+5^n\mathbb Z_5)
   =\alpha_m^{-n}I_n(a)
       -5^w\alpha_m^{-n-1}I_{n-1}(a).
\]
Indeed the stabilization is
$g_m-(5^m/\alpha_m)g_m(5z)$, and change of variable in the
integral introduces $1/5$. This also proves the distribution
identity. The program independently checks it on all twenty
unit residue classes modulo $25$:
\[
 \sum_{b=0}^4I_2(a+5b)
       =a_5(g_m)I_1(a)-5^w I_0(a)\quad(1\le a\le4).
\]
All polynomial path coordinates are integral, so these
measures are bounded in the displayed $\mathbb Z_5$ lattices.

For odd $m$ the twisted calculation stays at level $32$.
The Gauss identity and a change of variable give
\[
\begin{split}
24^w\tau(\chi_{24})I_{n,6}(a)
  =\sum_{v\bmod24}\chi_{24}(v)
    \int_{a/5^n+v/24}^{i\infty}
       g_m(z)\bigl(24\cdot5^nz-24a-5^nv\bigr)^w\,dz.
\end{split}
\]
This is the exact formula executed by the program. The
polynomials have integer coefficients; both $24^w$ and
$\tau(\chi_{24})$ are $5$-adic units, since
$\tau(\chi_{24})^2=24$. Thus this normalization changes
neither boundedness nor nonvanishing. No large-level
modular-symbol approximation is being substituted for the twist.

Let $\kappa$ denote the cyclotomic character on
$G_\infty=\operatorname{Gal}(\mathbb Q(\mu_{5^\infty})/\mathbb Q)$.
The Fourier and Gauss calculation of F.17, now in degree $w$,
identifies the Mellin transform of $\mu_m$ with Kato's
$L$-function evaluated at $\kappa^m$ times its argument.
It uses the critical value $r=m=k-1$ in
\citet[Theorem~16.2]{FA_Kato2004}. Therefore the required
finite-character identities identify the entire transform:
both sides are bounded Iwasawa measures, and a nonzero
bounded Iwasawa function has only finitely many zeros on
the open character disk. The required
specialization for $W_{m,d}(-t)$ is the scalar represented by
\[
 J_{m,d,t}=\int_{\mathbb Z_5^\times}x^{-m-t}\,d\mu_{m,d}(x),
 \qquad
 J_{m,d,t}\sim L_{p,\alpha_m}(g_{m,d})(\kappa^{-t}),
\]
where $\sim$ means equality up to a fixed nonzero period
scalar on the indicated parity component.

Here is a completely finite precision certificate. For units
$x\equiv a\pmod {25}$, one has
$x^{-m-t}\equiv a^{-m-t}\pmod {25}$. Since the measure is
integral and bounded, integrating the omitted variation
contributes an element of $25\mathbb Z_5^2$. Hence
\[
 J_{m,d,t}\equiv
 \sum_{\substack{1\le a<25\\5\nmid a}}
     a^{-m-t}\mu_{m,d}(a+25\mathbb Z_5)\pmod {25}.
\]
The unit root is lifted exactly modulo $25$ by the polynomial
above, so every term of this finite sum is exact modulo $25$.
For the twisted entries the displayed vector includes the
fixed unit factor $24^w\tau(\chi_{24})$.
\[
\begin{array}{cc|cc}
m&t&E_1&E_6\\\hline
1&0&(22,3)&(19,6)\\
1&1&(0,5)&(0,5)\\
1&2&(19,6)&(24,1)\\
2&0&(9,9)&(9,9)\\
2&1&(6,19)&(6,19)\\
2&2&(7,7)&(7,7)\\
3&0&(20,15)&(4,18)\\
3&1&(16,0)&(7,0)\\
4&0&(5,0)&(5,0)\\
4&1&(0,20)&(0,20)\\
5&0&(6,3)&(12,6)\\
6&0&(12,12)&(12,12).
\end{array}
\]
Every entry is nonzero. In particular, entries divisible by
$5$ have not been discarded or incorrectly called units;
their nonzero residue modulo $25$ is sufficient.

The two coordinates do not conceal a sum of unrelated
periods. On polynomial homology the involution used by the
program sends $z$ to $-\overline z$ and $P(X,Y)$ to
$P(-X,Y)$. It sends $I_n(a)$ to $(-1)^wI_n(-a)$.
For $r=m+t$, the character $x^{-r}$ has parity $(-1)^r$,
so $J_{m,d,t}$ belongs to the single eigenspace with sign
\[
 (-1)^{w+r}=(-1)^{t+1}.
\]
Both eigenspaces are one-dimensional. The program constructs
the conjugation matrix $C$ on the entire rational quotient,
checks $C^2=1$ and preservation of the period-functional
equations, and checks the displayed vector's eigenvalue
modulo $25$. The coefficient duality between polynomial
homology and the cohomological convention includes the
determinant twist of degree $w$, explicitly
\[
 (\operatorname{Sym}^w H)^\vee
     \simeq\operatorname{Sym}^w H\otimes\det(H)^{-w}
\]
for the two-dimensional coefficient space $H$.
It changes this sign by
$(-1)^w$, giving the parity $(-1)^r$ in the top-critical
interpolation formula. Thus each nonzero vector certifies
the corresponding single scalar $p$-adic $L$-value.

We give the required localization argument, since a nonzero
specialization does not say that its whole Iwasawa branch
is a unit. Fix one of the pairs $(m,t)$, and let
$T\subset W_{m,d}$ be a stable lattice. Write
\[
 \Lambda=\mathbb Z_5[[G_\infty]],\qquad
 \mathcal S=\operatorname{Sel}_\infty(T),\qquad
 X=\mathcal S^\vee.
\]
By \citet[\S17.3 and Theorem~17.4(1)]{FA_Kato2004},
$X$ is a finitely generated torsion $\Lambda$-module.
Let $\mathfrak p_t$ be the height-one prime on the appropriate
$\Delta$ component defined, after inverting $5$, by evaluation
at $\kappa^{-t}$. The just-proved nonzero value and
Theorem~17.4(2) imply
\[
 X_{\mathfrak p_t}=0.
\]
Consequently
\[
 \bigl(X\otimes_{\Lambda,\kappa^{-t}}\mathbb Z_5\bigr)
          \otimes\mathbb Q_5=0.
\]
The integral tensor product is finitely generated over
$\mathbb Z_5$, hence finite. Pontryagin duality identifies it
exactly with
\[
 \bigl(\mathcal S(-t)^{G_\infty}\bigr)^\vee.
\]
Indeed the contragredient action on the dual changes the
twist $-t$ into $t$, and taking coinvariants imposes
$g=\kappa(g)^{-t}$ on $X$. This verifies both the character
and the sign of the specialization, and proves that the
desired eigenspace of the discrete Selmer group has bounded
exponent. It makes no assertion about the other eigenspaces.

Finally put $M=W_{m,d}(-t)$ and $B=M/T(-t)$.
The argument of F.17 gives
$W_{m,d}^{G_{\mathbb Q(\mu_{5^\infty})}}=0$ for every
$m\ge1$: its two CM characters have distinct Hodge--Tate
weights, while complex conjugation exchanges them. A
nonzero invariant subspace would force this irreducible
representation to factor through an abelian group and
contradict its two conjugation eigenspaces. Twisting does
not change the restriction to the cyclotomic kernel.
Thus $B^{G_{\mathbb Q(\mu_{5^\infty})}}$ is finite and
inflation--restriction has bounded-exponent kernel.

All local finite groups for $M$ are zero, because it is a
direct summand of one of the dual layers in the table and
F.9 proves that assertion for the entire layer. Fix
$\Sigma=\{2,3,5,\infty\}$. A finite Selmer class is
unramified outside this set, and the finitely generated
group $H^1(G_{\mathbb Q,\Sigma},T(-t))$ supplies the integral
lattice in rational cohomology. If
$H_f^1(\mathbb Q,M)$ were nonzero, reduction by this lattice
would give a nonzero divisible $5$-primary subgroup of
$H^1(G_{\mathbb Q,\Sigma},B)$.
Its restriction is locally trivial and therefore belongs
to the ordinary kernel in Kato's Proposition~17.2, after
the stated twist; being a restriction of global classes,
it lies in $\mathcal S(-t)^{G_\infty}$.
That target has bounded exponent, so the divisible image
is zero. The remaining restriction kernel also has bounded
exponent, and cannot contain the original nonzero divisible
group. This contradiction proves
\[
 H_f^1(\mathbb Q,W_{m,d}(-t))=0
        \qquad(m+2t\le6).
\]
All balanced summands were already handled, so every defect
in the statement vanishes. The exact global tables now
follow from the successive obstruction and fibre argument
proved in F.9.
\end{proof}

\subsection{Integral real metabelian sections and Fermat descent}

\begin{statement}{F.19}{Proved}{The complete real pro-two metabelian section set}
Let $X=\mathbb P^1_{\mathbb R}\setminus\{0,1,\infty\}$,
with the tangential base point at $0$ pointing towards $1$.
Let $P$ be its geometric free pro-$2$ fundamental group,
$M=P/\overline{P''}$, and $N_2=P/\overline{\gamma_3P}$.
Choose the real-path generators $x,y$ around $0,1$, and
write $\sigma$ for complex conjugation. Then
\[
 H^1_{\rm cts}(G_{\mathbb R},M)
 \ \xrightarrow{\ \sim\ }\
 H^1_{\rm cts}(G_{\mathbb R},N_2)
 =\{[1],[x],[y]\}.
\]
Their images in
$H^1(G_{\mathbb R},\mathbb Z_2(1)^2)=(\mathbb Z/2)^2$
are $(0,0),(1,0),(0,1)$ respectively. They are the
sections of the three real connected components
$(0,1),(-\infty,0),(1,\infty)$.

For $r\ge1$ let
\[
 D_r=(\mathbb Z/2^r)^2,\qquad
 K_r=\ker(M\longrightarrow D_r).
\]
These are the kernels induced in the base's metabelian
quotient by the Fermat covers
$u^{2^r}=z$, $v^{2^r}=1-z$.
Exactly $[1]$ among the three base classes descends to
$K_r\rtimes G_{\mathbb R}$, and
\[
 |H^1(G_{\mathbb R},K_r)|=4.
\]
Representatives are the coboundaries
$\sigma\mapsto h^{-1}\sigma(h)$ for
$h=x^{2^{r-1}i}y^{2^{r-1}j}$, $i,j\in\{0,1\}$.
All four classes at level $r+1$ map to the distinguished
class at level $r$. The groups $K_r$ are the induced
relative kernels; this assertion does not identify them
with the absolute metabelian completions of the covers.
\end{statement}
\begin{proof}
Real paths give $\sigma(x)=x^{-1}$ and
$\sigma(y)=y^{-1}$. These are the explicit generators
in \citet[Example~5.6]{FA_Wickelgren2014}; the real path
from her base point $1/2$ to our tangent identifies the
two actions. Her Theorem~1.1 also identifies the image
of the class-two section set in the abelian section set.
We compute the entire integral fibres directly.

Use the commutator $c=[x,y]=xyx^{-1}y^{-1}$.
The unique normal forms in $N_2$ are $x^a y^b c^d$,
$(a,b,d)\in\mathbb Z_2^3$, with multiplication and action
\[
 (a,b,d)(a',b',d')=(a+a',b+b',d+d'-a'b),
 \qquad \sigma(a,b,d)=(-a,-b,d).
\]
A cocycle is determined by its value $g$ at $\sigma$;
the cocycle condition is exactly
\[
             g\sigma(g)=1
       \quad\Longleftrightarrow\quad 2d+ab=0.
\]
Thus $ab$ must be even and $d=-ab/2$ is uniquely
determined. Twisted conjugation by $h=(u,v,w)$ gives
\[
 (a,b,d)\longmapsto
 (a-2u,b-2v,d+av+ub-2uv).
\]
Choose $u=(a-\epsilon)/2$, $v=(b-\eta)/2$, where
$\epsilon,\eta\in\{0,1\}$ are the residues of $a,b$.
The condition excludes $(\epsilon,\eta)=(1,1)$ and
reduces $g$ to exactly one of $1,x,y$. Their abelian
residues distinguish them. Kummer signs of $(z,1-z)$
give the three connected components asserted above.
The omitted fourth abelian class has a nonzero integral
obstruction, even though rationalization would remove it.

For the full metabelian group put
$A=M^{\rm ab}=\mathbb Z_2^2$ and
$\Lambda=\mathbb Z_2[[A]]$; denote conjugation by $x,y$
on $M'$ by multiplication by $X,Y$.
The completed Alexander calculation is
\[
                 M'=\Lambda c\simeq\Lambda.
\]
This is the rank-two case of
\citet[Proposition~8 and \S5]{FA_Segal2023}.
For clarity, the completed Fox boundary sends the two
basis vectors to $X-1,Y-1$. Its kernel is generated
freely by $(-(Y-1),X-1)$: the two parameters are prime
and relatively prime in $\mathbb Z_2[[X-1,Y-1]]$.
The commutator gives this generator up to a unit.

Let $\overline{\phantom a}$ denote inversion on $A$,
extended linearly and continuously to $\Lambda$.
The word identity
\[
 \sigma(c)=[x^{-1},y^{-1}]
              =(xy)^{-1}c(xy)
\]
shows that the twisted involutions on $M'$ belonging
to the three cocycle lifts $s=1,x,y$ are respectively
\[
 \tau_1(r)=X^{-1}Y^{-1}\overline r,\qquad
 \tau_x(r)=Y^{-1}\overline r,\qquad
 \tau_y(r)=X^{-1}\overline r.
\]
On the point masses in the completed group module
$\Lambda=\varprojlim_n(\mathbb Z/2^n)[A/2^nA]$
these are the affine
involutions $v\mapsto-v+q$, for
$q=(-1,-1),(0,-1),(-1,0)$. In each case a coordinate
of $q$ is odd. Partition $A$ by the parity of that
coordinate. The two resulting clopen pieces are
exchanged by $\tau$, so the completed group module
decomposes topologically as $C\oplus\tau C$.
If $r+\tau r=0$, its component $r_0\in C$ satisfies
$r=r_0-\tau r_0=(\tau-1)(-r_0)$. Consequently
\[
 H^1(G_{\mathbb R},(M')^{s})
       =\ker(1+\tau_s)/\operatorname{im}(\tau_s-1)=0
                  \qquad(s=1,x,y).
\]
Every metabelian cocycle first has one of the three
allowed abelian residues, by its class-two image.
Twisted conjugation makes its abelian value equal
to that of the indicated $s$ exactly. Its remaining
factor is a cocycle in the displayed twisted module,
and the explicit coboundary just constructed removes
it. This proves both surjectivity and injectivity of
the first map, without an inverse-limit ambiguity.

For Fermat descent the action on $D_r$ is inversion,
so $H^1(G_{\mathbb R},D_r)=D_r/2D_r$.
The three base classes have the same three distinct
residues there; only $[1]$ maps to zero.
The nonabelian exact sequence identifies the fibre
over $[1]$ with
\[
 D_r^{\sigma}/\operatorname{im}(M^{\sigma})
   =D_r[2]=(\mathbb Z/2)^2.
\]
Indeed an element of $M^{\sigma}$ has an abelian
image $a=-a$ in the torsion-free group $A$, hence
has zero image in $D_r$. The connecting cocycles
are precisely those in the statement. At the next
level their conjugating elements already lie in
$K_r$, so their images become coboundaries in
$K_r$. Geometrically the Fermat curve has four real
arcs, distinguished by the signs of $u,v$; squaring
both coordinates maps all of them to the positive
arc. This matches the calculated transition map.

The exact program enumerates every class-two
cocycle and its full twisted-conjugacy orbit over
$\mathbb Z/2^n$ for $1\le n\le8$, and every basis
element of the three Alexander permutation modules.
The four finite classes have representatives
$1,c^{2^{n-1}},x,y$, with orbit sizes
$2^{2n-2},2^{2n-2},2^{2n-1},2^{2n-1}$.
Reduction sends their class indices to $(0,0,2,3)$.
Thus the extra finite central class disappears in
the inverse limit, as the integral proof requires.
The program also verifies every displayed Fermat
representative through $r=8$ with exact integer
coordinates. These are finite checks of the stated
formulas; the module argument proves all levels.
\end{proof}

\subsection{A complete integral global section calculation at class two}

\begin{statement}{F.20}{Proved}{Global pro-five class-two sections as explicit tame congruences}
Let $S$ be a finite set of rational primes containing $5$,
$G_S=G_{\mathbb Q,S\cup\{\infty\}}$, and
$R_S=\mathbb Z[S^{-1}]$. Let $U$ be the class-two
quotient of the geometric pro-$5$ fundamental group
of $\mathbb P^1\setminus\{0,1,\infty\}$, with its
actual arithmetic action and tangential section at $0+1$.
No crystalline condition is imposed. Identify
\[
 H^1(G_S,U^{\rm ab})
    =\mathbb Z_5^S\times\mathbb Z_5^S
\]
by the Kummer classes of the prime units in $S$.
For each $q\in S$ with $e_q=v_5(q-1)>0$, choose a
primitive root $g_q$ modulo $q$ and put
\[
 \ell_q(\alpha)=\sum_{p\in S\setminus\{q\}}
             \alpha_p\log_{g_q}(p\bmod q)
                  \pmod {5^{e_q}},\qquad
 \Omega_q(\alpha,\beta)
       =\ell_q(\alpha)\beta_q-\alpha_q\ell_q(\beta).
\]
Then abelianization is a bijection
\[
 H^1_{\rm cts}(G_S,U)\ \xrightarrow{\ \sim\ }\
 \bigl\{(\alpha,\beta)\in\mathbb Z_5^S\times\mathbb Z_5^S:
          \Omega_q(\alpha,\beta)=0
              \text{ for every such }q\bigr\}.
\]
In particular every nonempty integral fibre has exactly
one conjugacy class. For the two explicit sets this reads
\[
 \begin{array}{c|c}
 S&H^1(G_S,U)\text{ in Kummer coordinates}\\\hline
 \{2,3,5\}&\mathbb Z_5^6\\
 \{2,3,5,11\}&
 \{(\alpha,\beta)\in\mathbb Z_5^8:
                         \Omega_{11}(\alpha,\beta)=0\}
 \end{array}
\]
where the second equation is modulo $5$ and
\[
 \Omega_{11}=(\alpha_2+3\alpha_3+4\alpha_5)\beta_{11}
          -\alpha_{11}(\beta_2+3\beta_3+4\beta_5).
\]
For the second set there are exactly $90\,625$
admissible abelian pairs modulo $5$, out of $390\,625$.
The pair of units $(2,11)$ is obstructed; $(-10,11)$
is unobstructed and is the pair of the point $z=-10$.

Let $U_r$ be the Heisenberg quotient with all three
integral coordinates reduced modulo $5^r$. For
$S=\{2,3,5\}$ its section set has $5^{6r}$ classes,
one over each abelian class. For
$S=\{2,3,5,11\}$ it has
\[
       5\cdot90\,625\cdot5^{8(r-1)}
\]
classes; each admissible abelian class has exactly
five lifts. These finite extra lifts do not create
additional integral fibres.

For the induced Fermat kernel
$K_r=\ker(U\to(\mathbb Z/5^r(1))^2)$,
$H^1(G_S,K_r)\to H^1(G_S,U)$ is injective, with
image the displayed congruence locus intersected
with $5^r\mathbb Z_5^{2|S|}$.
For the two explicit sets the image is exactly
$5^r\mathbb Z_5^6$ and $5^r\mathbb Z_5^8$,
respectively. Only the distinguished base class
descends to every level of this induced Fermat
$5$-power system. These statements compute the
specified global class-two section sets; they do
not compute the unrestricted global metabelian
section set in F.12.
\end{statement}
\begin{proof}
With the stated tangential base point the arithmetic
action on the free generators has the form
\[
 g(x)=x^{\chi_5(g)},\qquad
 g(y)=f_g^{-1}y^{\chi_5(g)}f_g,\qquad f_g\in[P,P].
\]
See \citet[(9), \S3.11]{FA_Wickelgren2012}.
The conjugator is central at class two and disappears
from the action there. Since $2$ is a unit in
$\mathbb Z_5$, integral logarithmic coordinates give
\[
 U=\mathbb Z_5(1)^2\times\mathbb Z_5(2),\qquad
 (a,b,t)(a',b',t')=
 (a+a',b+b',t+t'+\tfrac12(ab'-ba')),
\]
and $g(a,b,t)=(\chi_5(g)a,\chi_5(g)b,\chi_5(g)^2t)$.
Thus the action factors through $G_S$ and the central
extension has kernel $T=\mathbb Z_5(2)$.

We first evaluate its integral cohomology, including
its torsion. \citet[Theorem~4.5]{FA_Levine1989}
identifies $H^1(\mathbb Q,\mathbb Z_5(2))$ with
$K_3(\mathbb Q)^{\rm ind}\otimes\mathbb Z_5$.
The latter is zero: $K_3(\mathbb Q)=\mathbb Z/48$,
as recalled and recovered in that paper's introduction.
Inflation from $G_S$ is injective, hence
\[
                     H^1(G_S,T)=0.
\]
The integral comparison of
\citet[Theorem~5.1]{FA_NguyenQuangDo1986} is
\[
 H^2(G_S,T)\simeq K_2(R_S)\{5\}.
\]
Here the notation on the right means the $5$-primary
part of the finite group, not its rationalization.
Localization and the computations
$K_2(\mathbb Z)=\mathbb Z/2$ and
$K_2(\mathbb F_q)=0$ give
\[
 0\longrightarrow K_2(\mathbb Z)
 \longrightarrow K_2(R_S)
 \xrightarrow{(\partial_q)_{q\in S}}
        \bigoplus_{q\in S}\mathbb F_q^\times
 \longrightarrow0.
\]
The map on $K_1$ following this sequence is zero
because $\mathbb Z^\times\to R_S^\times$ is injective.
The localization and low $K$-group calculations are
given in \citet[III.5.2.2, III.6.1.1 and V.6.6]{FA_Weibel2013}.
Taking the $5$-primary part therefore proves the explicit
formula
\[
        H^2(G_S,T)\simeq
        C_S:=\bigoplus_{\substack{q\in S\\e_q>0}}
                         \mathbb Z/5^{e_q}.
\]
For $S=\{2,3,5\}$ this is zero; adjoining $11$ gives
one cyclic group of order $5$.

The Kummer sequence and $\operatorname{Pic}(R_S)=0$
give $H^1(G_S,\mathbb Z_5(1))=\mathbb Z_5^S$:
$R_S^\times=\{\pm1\}\times\langle p:p\in S\rangle$,
and the sign contributes nothing at $5$.
For cocycles $a,b$ representing $\alpha,\beta$, the
central cocycle equation in the displayed group law
has obstruction
\[
 \tfrac12\bigl([a\smile b]-[b\smile a]\bigr)
                       =[a\smile b]\quad\text{in }H^2(G_S,T).
\]
Graded commutativity gives the equality. The Chern
class sends the Steinberg symbol $\{u,v\}$ to
$-\kappa(u)\smile\kappa(v)$; see
\citet[\S2.2.4.4 and Lemma~2.2.4.4]{FA_Soule1978}.
Under the tame maps,
\[
 \partial_q\{u,v\}
       =(-1)^{v_q(u)v_q(v)}
          \frac{u^{v_q(v)}}{v^{v_q(u)}}\pmod q.
\]
Its $5$-primary logarithm is precisely $\Omega_q$;
the sign factor has order at most two and disappears.
The fixed global sign from the Chern convention
does not affect the zero condition. First this
computes the obstruction for integral prime
exponents; continuity and $\mathbb Z_5$-bilinearity
extend it to all the displayed coordinates.
Thus the obstruction vanishes exactly on the stated
congruence locus. Every nonempty fibre of the
central nonabelian cohomology map is a quotient
of a torsor under $H^1(G_S,T)=0$, so is a singleton.
In particular the result is an evaluated section
set, not just a reduction to an unnamed cup product.

For the numerical instance choose $g_{11}=2$.
The discrete logarithms of $2,3,5$ are $1,8,4$
modulo $10$, and reduce to $1,3,4$ modulo $5$.
The resulting alternating matrix on $\mathbb F_5^4$
has rank two. The map
\[
 \alpha\longmapsto
   (\alpha_2+3\alpha_3+4\alpha_5,\alpha_{11})
\]
has a two-dimensional kernel. There are
$5^2+(5^2-1)5=145$ ordered dependent pairs in
$\mathbb F_5^2$, each with $5^4$ inverse images.
Hence the number is $145\cdot625=90\,625$.
The exact program independently enumerates all
$5^8=390\,625$ ordered pairs and computes each
tame symbol multiplicatively in $\mathbb F_{11}^\times$.
Its five obstruction counts are
$(90\,625,75\,000,75\,000,75\,000,75\,000)$.
For $(-10,11)$ the first logarithm is $1+4=0$,
whereas for $(2,11)$ it is $1$.

We also evaluate the finite fibres, to control the
integral limit. Put $T_r=T/5^rT$.
For finite $S$, cohomology with finite $5$-primary
coefficients is finite in every degree and vanishes
above degree two, by
\citet[I, Theorem~4.10(c) and Corollary~4.15]{FA_Milne2006ADT}.
The real-place terms in the former theorem vanish
because their Galois groups have order two.
Consequently
$H^3(G_S,T)=0$: the possible derived-limit term
is zero for an inverse system of finite groups.
The coefficient sequence gives
\[
 H^1(G_S,T_r)=C_S[5^r],\qquad
 H^2(G_S,T_r)=C_S/5^rC_S.
\]
Also $H^0(G_S,\mu_{5^r})=0$, so every nonempty
finite fibre is a torsor under $C_S[5^r]$, without
a further stabilizer quotient. Its obstruction is
the same $\Omega_q$ reduced modulo
$5^{\min(e_q,r)}$.
For the first example this proves $5^{6r}$ classes.
For the second, admissibility depends only on
the eight coordinates modulo $5$, giving
$90\,625\cdot5^{8(r-1)}$ abelian classes and
five lifts of each. Under reduction $r+1\to r$
the map on the central $H^1$ groups is multiplication
by $5$ under their Bockstein identification with
subgroups of $C_S$. This follows from the morphism
of coefficient sequences whose left map on $T$
is multiplication by $5$. In particular it is
zero in the second example. For a fixed integral
abelian cocycle these finite torsors therefore
create no extra inverse-limit fibre, consistently
with the direct integral calculation above.

Finally $D_r=(\mathbb Z/5^r(1))^2$ has no
$G_S$-invariants. The nonabelian exact sequence
for $1\to K_r\to U\to D_r\to1$ makes
$H^1(G_S,K_r)\to H^1(G_S,U)$ injective, with
image the kernel of the finite Kummer map.
In our coordinates that map is reduction modulo
$5^r$. This proves the descent assertion. In
the second example all the obstruction terms
are divisible by $5^{2r}$ on this sublattice,
so its modulo-$5$ equation is automatic.
The intersection of the sublattices over all
$r$ is zero. As in F.19 these are induced
relative kernels; F.11 does not permit their
replacement by absolute class-two groups of
the covers. The order-four quotient $q_2$
does not factor through a pro-$5$ group at all.

This also computes an additional cusp collision:
the tangential sections at $0+1$, $1-1$, and
the image of $0+1$ under $z\mapsto1/z$ all have
zero $5$-adic abelian Kummer class, by
\citet[Lemma~3.16]{FA_Wickelgren2012}, and hence
the same integral class-two section class.
Neither this collision nor the explicit finite
ramification computation evaluates the higher
global metabelian fibres left in F.12.
\end{proof}

\subsection{Unrestricted global class two and the derived-limit obstruction}

\begin{statement}{F.21}{Proved}{The unrestricted integral class-two section set}
Let $\mathcal P$ be the set of rational primes and put
\[
 \mathcal C(J)=\bigl\{(a_j)\in\prod_{j\in J}\mathbb Z_5:
      \text{for every }n,\ a_j\in5^n\mathbb Z_5
                         \text{ for all but finitely many }j\bigr\}.
\]
Let $I=\{q\in\mathcal P:5\mid q-1\}$ and
$e_q=v_5(q-1)$. For each $q\in I$ choose a primitive
root $g_q$ modulo $q$, and choose integral lifts
\[
 L_q(p)\in\mathbb Z_5,\qquad
 L_q(p)\equiv\log_{g_q}(p\bmod q)\pmod {5^{e_q}}
                 \quad(p\ne q).
\]
For $\alpha,\beta\in\mathcal C(\mathcal P)$ define
\[
 T_q(\alpha,\beta)=
 \beta_q\sum_{p\ne q}\alpha_pL_q(p)
       -\alpha_q\sum_{p\ne q}\beta_pL_q(p).
\]
All sums converge, and $T(\alpha,\beta)\in\mathcal C(I)$.
For the actual pro-$5$ class-two tripod group $U$ of F.20,
abelianization is a bijection
\[
 H^1_{\rm cts}(G_{\mathbb Q},U)\ \xrightarrow{\ \sim\ }\
 \left\{(\alpha,\beta)\in\mathcal C(\mathcal P)^2:
       \left(\frac{T_q(\alpha,\beta)}{5^{e_q}}\right)_{q\in I}
                                      \in\mathcal C(I)\right\}.
\]
Membership includes integrality of every displayed quotient.
Thus the finite tame congruences must be supplemented by
the condition that the divided coordinates tend to zero
$5$-adically. The criterion is independent of all the
choices of logarithm lifts and primitive roots.

For the induced relative Fermat kernel
$K_r=\ker(U\to(\mathbb Z/5^r(1))^2)$, its section
set injects into this locus as its intersection with
$5^r\mathcal C(\mathcal P)^2$. Only the distinguished
base class descends to every $r$.

More precisely, as an algebraic $\mathbb Z_5$-module,
\[
 H^2_{\rm cts}(G_{\mathbb Q},\mathbb Z_5(2))
   \simeq\mathcal C(I)/D\mathcal C(I),\qquad
                 D(a)_q=5^{e_q}a_q.
\]
The denominator is its actual image, without taking a
closure. The kernel of the finite-coefficient projections is
\[
 \frac{\{a\in\prod_{q\in I}\mathbb Z_5:Da\in\mathcal C(I)\}}
      {\mathcal C(I)}
 \simeq
 \varprojlim{}^1_n
       H^1(G_{\mathbb Q},\mathbb Z/5^n(2)).
\]
This kernel occurs for actual Kummer cup products: there
are abelian pairs lifting to $U/5^n$ for every $n$ but
admitting no integral class-two lift. Here $U/5^n$ means
the finite Heisenberg quotient of F.20.
\end{statement}
\begin{proof}
\emph{The continuous cochain comparison.}
For $j=1,2$ put
\[
 C_{0,j}^{\bullet}=\varinjlim_{S\ni5,\ S\ {\rm finite}}
           C^{\bullet}_{\rm cts}(G_S,\mathbb Z_5(j)),
\]
using inflation. These are complexes of torsion-free
$\mathbb Z_5$-modules. In every cochain degree,
\[
 C_{0,j}^{\bullet}/5^n
       =C^{\bullet}_{\rm cts}(G_{\mathbb Q},\mathbb Z/5^n(j)).
                                                        \tag{F.21.1}
\]
Indeed a continuous cochain with finite values factors
through a finite Galois quotient, whose ramification set
is finite. Conversely such a cochain lifts to a continuous
$\mathbb Z_5$-valued cochain by choosing lifts of its
finitely many values. Its lift still factors through the
same quotient as a function. The coefficient action factors
through $G_S$ after adjoining $5$ to $S$.
The inverse limit of (F.21.1) is therefore exactly the
continuous cochain complex with coefficients $\mathbb Z_5(j)$.

By F.20 and the finite-coefficient vanishing used there,
\[
 H^i(C_{0,2}^{\bullet})=
 \begin{cases}
 A=\displaystyle\bigoplus_{q\in I}\mathbb Z/5^{e_q},&i=2,\\
 0,&i\ne2.
 \end{cases}                                           \tag{F.21.2}
\]
Here filtered colimits are exact. In degree zero the
cyclotomic square has no invariants; in degree one use
F.20; above degree two use
\citet[I, Theorem~4.10(c) and Corollary~4.15]{FA_Milne2006ADT}
and the finite-group inverse-limit argument at each fixed $S$.
The comparison with the tame direct sum is compatible
with enlarging $S$.

Let $P_1=P_0=\bigoplus_{q\in I}\mathbb Z_5$, and place
the diagonal complex $P_1\xrightarrow{D}P_0$ in degrees
one and two. Choose a degree-two cocycle $z_q$ representing
the $q$th generator of $A$. Since its order is $5^{e_q}$,
choose a cochain $h_q$ in $C_{0,2}^1$ with
$dh_q=5^{e_q}z_q$. These choices define a quasi-isomorphism
\[
       [P_1\xrightarrow{D}P_0]\longrightarrow C_{0,2}^{\bullet}.
                                                        \tag{F.21.3}
\]
It remains a quasi-isomorphism modulo $5^n$: multiplication
by $5^n$ is injective in every term, so this follows from
the two coefficient exact sequences. All reduction maps
on cochain groups, and on the terms $P_i/5^n$, are
surjective. Their inverse limits preserve this
quasi-isomorphism. One can see this directly from the
degreewise short exact sequence with map
$1-\mathrm{shift}$ on the product of the complexes;
the product of quasi-isomorphisms is a quasi-isomorphism.
Thus (F.21.3) completes to the exact cochain model
\[
 [\mathcal C(I)\xrightarrow{D}\mathcal C(I)]
          \simeq C^{\bullet}_{\rm cts}
                         (G_{\mathbb Q},\mathbb Z_5(2)). \tag{F.21.4}
\]
In particular its degree-one cohomology is zero and its
degree-two cohomology is the stated cokernel. This proves
the derived-completion calculation without imposing a
Mittag--Leffler hypothesis on the global cohomology groups.

\emph{The cup product, including the nonseparated part.}
Kummer theory identifies
$H^1(G_{\mathbb Q},\mathbb Z_5(1))=\mathcal C(\mathcal P)$.
To identify its cup product in (F.21.4), it is essential
to work with cochains: equality on finite-support inputs
alone would not determine a map into this potentially
non-Hausdorff cohomology quotient.

Put $B=\bigoplus_{p\in\mathcal P}\mathbb Z_5$. For every
prime $p$, choose compatible $5$-power roots of $p$ and
the associated integral Kummer cocycle. It is unramified
outside $\{5,p\}$ and gives a chain map
$B[-1]\to C_{0,1}^{\bullet}$. Cup product is then a
chain map
\[
                 (B\otimes B)[-2]\longrightarrow C_{0,2}^{\bullet}.
\]
The free module $B\otimes B$ is projective. Its map in
degree-two cohomology is the tame-symbol map proved in
F.20; absorb the fixed Chern-class sign into the chosen
identification of $A$. Lift this map to $P_0$ using the
integers $L_q(p)$. Each column has support in at most
the two primes occurring in that column. The original
cup cochain and the cochain obtained from this lifted
map through (F.21.3) have the same cohomology class on
every basis vector. Choose their integral bounding
cochains on these basis vectors. Since the source is
free, these choices give an actual $\mathbb Z_5$-linear
chain homotopy.

Every map and this homotopy extend to the $5$-adic
completions: a $\mathbb Z_5$-linear map preserves the
$5^n$ filtrations. Also
$\widehat{B\otimes B}=\mathcal C(\mathcal P\times\mathcal P)$,
and $(\alpha_p\beta_r)_{p,r}$ belongs to this module.
Cup product on continuous cochains is continuous for
these filtrations. Consequently the completed chain
homotopy proves that the actual integral cup is precisely
the class of the displayed vector $T(\alpha,\beta)$
in (F.21.4). This is an equality of cohomology classes
proved before taking the quotient, including its
nonseparated part.

For each fixed $q$ the sums defining $T_q$ converge,
since the coefficient sequences tend to zero and all
$L_q(p)$ are integral. Their absolute values are bounded
uniformly in $q$. The factors $\alpha_q,\beta_q$ then
show that $T\in\mathcal C(I)$. Changing a logarithm lift
changes the lifted map $B\otimes B\to P_0$ by $D$ times
an integral map to $P_1$ with finite columns. Its
completion changes $T$ by $D\mathcal C(I)$.
Changing a primitive root multiplies the corresponding
coordinate by a $5$-adic unit, together with such a
change of lift, so it does not alter the criterion.
The central cocycle equation of F.20 now proves
existence exactly when $D^{-1}T\in\mathcal C(I)$.
Its integral fibres are singletons because
$H^1(G_{\mathbb Q},\mathbb Z_5(2))=0$.
For the Fermat assertion the finite cyclotomic quotient
has no Galois invariants, exactly as in F.20; its
cohomology map is coordinatewise reduction modulo
$5^r$. The intersection of these submodules over
all $r$ is zero.

\emph{The derived-limit term explicitly.}
Reducing (F.21.3) modulo $5^n$ gives
\[
 H^1(G_{\mathbb Q},\mathbb Z/5^n(2))=A[5^n],\qquad
 H^2(G_{\mathbb Q},\mathbb Z/5^n(2))=A/5^nA.
\]
The transition on the first groups is multiplication by
$5$. In the inverse limit of the second groups, choose
for each coordinate its representative in
$\{0,\ldots,5^{e_q}-1\}$. The inverse-limit condition
says exactly that these representatives belong to
$\mathcal C(I)$. Hence the projection of the cokernel
in (F.21.4) onto $\varprojlim A/5^nA$ is surjective,
and its kernel is the stated quotient of divided
sequences. The cochain product sequence with
$1-\mathrm{shift}$ identifies this kernel with the
displayed $\varprojlim{}^1$. No closure of
$D\mathcal C(I)$ is taken in this argument.

\emph{An actual cup invisible at every finite coefficient level.}
Choose distinct primes $q_n$ such that
\[
 v_5(q_n-1)=n,\qquad
                 5^{(q_n-1)/5}\not\equiv1\pmod {q_n}.
                                                        \tag{F.21.5}
\]
Such primes exist for every $n$. Indeed
$5^{1/5}\notin E_n=\mathbb Q(\zeta_{5^{n+1}})$:
the non-Galois field $\mathbb Q(5^{1/5})$ cannot be a
subfield of an abelian Galois extension of $\mathbb Q$.
Since $E_n$ contains $\mu_5$, Kummer theory then gives
$[E_n(5^{1/5}):E_n]=5$. The cyclotomic automorphism
with exponent $1+5^n$ extends by fixing the chosen
root of $X^5-5$; compose this extension with multiplication
of that root by $\zeta_5$. Thus the Galois group over
$\mathbb Q$ contains the automorphism
\[
 \zeta_{5^{n+1}}\longmapsto
        \zeta_{5^{n+1}}^{1+5^n},\qquad
          5^{1/5}\longmapsto\zeta_5\,5^{1/5}.
\]
Its Frobenius conjugacy class has positive density by
the Chebotarev theorem; see
\citet[Theorem~8.31]{FA_Milne2020ANT}.
The first action gives valuation exactly $n$, and the
second gives the nontrivial fifth-power residue in
(F.21.5). Both properties are unchanged under conjugacy.

Set $\alpha_5=1$ and $\alpha_p=0$ otherwise; set
$\beta_{q_n}=5^n$ and set all its other coordinates
to zero. These are valid elements of
$\mathcal C(\mathcal P)$. The only nonzero tame
coordinates are
\[
 T_{q_n}=5^nL_{q_n}(5),\qquad
       T_{q_n}/5^{e_{q_n}}=L_{q_n}(5)\in\mathbb Z_5^\times.
\]
They satisfy every finite tame congruence, but their
divided coordinates do not tend to zero. Thus the
integral cup is nonzero in the derived-limit kernel.
Every finite quotient has a lift, as follows either
from its zero finite obstruction or from F.20 applied
to the finitely many coordinates surviving modulo
$5^r$. The section classes of these finite lifts
cannot be chosen compatibly. This does not contradict
F.15, which concerns a compatible inverse-limit class.

The compatibility defect also has an explicit
$\varprojlim{}^1$ representative. Write $a_k$ for the
generator of the cyclic summand $\mathbb Z/5^k$ at
$q_k$, and put $L_k=L_{q_k}(5)$. In the diagonal
complex, a finite bounding cochain modulo $5^r$ is
$\sum_{k<r}L_k e_{q_k}$; successive reductions differ
by $L_r e_{q_r}$. Under the Bockstein identification
with $A[5^r]$ this gives, up to the fixed connecting-map
sign, the sequence $\gamma_r=L_r a_r$.
If it were zero in the derived limit, there would be
$b_r\in A[5^r]$ with $b_r-5b_{r+1}=\gamma_r$.
Iteration in the finite cyclic $q_k$-coordinate forces
\[
             b_r(q_k)=5^{k-r}L_k a_k\ne0
                         \qquad(k\ge r).
\]
For fixed $r$ this has infinite support, contrary to
the finite support required for $b_r\in A[5^r]\subset A$.
Thus the finite torsors have an explicitly nonzero
compatibility obstruction.

The exact program certifies (F.21.5) for $n=1,\ldots,8$,
using respectively
\[
 11,\ 101,\ 751,\ 13751,\ 37501,\ 62501,\
                        937501,\ 9375001.
\]
Primality is certified by complete factorization of
$q_n-1$ and an element of exact order $q_n-1$;
all power-residue and discrete-logarithm equalities
are checked exactly. The infinite existence argument
is Chebotarev, not extrapolation from this list.
The unrestricted class-two problem is therefore
evaluated, while the higher global metabelian groups
in F.12 remain a separate problem.
\end{proof}

\subsection{Integral global sections through class four}

\begin{statement}{F.22}{Proved}{The complete class-three and class-four sets over \(\mathbb Z[1/30]\)}
Put \(S=\{2,3,5\}\), \(G=G_{\mathbb Q,S\cup\{\infty\}}\), and let
\(P\) be the geometric free pro-\(5\) group of
\(\mathbb P^1-\{0,1,\infty\}\), with the arithmetic action defined
by the unit tangent \(0+\!1\). Write
\[
 N_n=P/\overline{\gamma_{n+1}P},\qquad
                 \mathcal S_n=H^1_{\rm cts}(G,N_n).
\]
These are integral section classes, with ramification restricted
to \(S\) and no additional crystalline condition. The actual
arithmetic action, including its Ihara conjugator, is used.
The successive maps have the following complete description:
\[
\begin{array}{c|c|c}
 n&\mathcal S_n\longrightarrow\mathcal S_{n-1}
       &\text{fibre over every class}\\ \hline
 2&\text{bijection}&\text{one point}\\
 3&\text{surjection}&\mathbb Z_5^2\text{-torsor}\\
 4&\text{surjection}&(\mathbb Z/5)^3\text{-torsor}.
\end{array}                                                   \tag{F.22.1}
\]
Here \(\mathcal S_1=\mathbb Z_5^6\), with the six Kummer
coordinates of \(2,3,5\) in the two abelian factors. In
particular, after choosing origins in the indicated torsors,
\[
 \mathcal S_3\simeq\mathbb Z_5^8,\qquad
 \mathcal S_4\simeq\mathbb Z_5^8\times(\mathbb Z/5)^3
                                                               \tag{F.22.2}
\]
as sets. No canonical coordinates on these torsors are asserted.

The \(125\) class-four lifts can be written explicitly.
Let \(e_1,e_2,e_3\) be an integral Hall basis of
\(\gamma_4P/\gamma_5P\), put
\[
       z_4(g)=\frac{\chi_5(g)^4-1}{5}\in\mathbb Z_5,
                                                               \tag{F.22.3}
\]
and choose one class-four cocycle \(c_g\) above a given
class-three cocycle. Its \(125\) distinct lift classes are
represented by
\[
             c_g\prod_{i=1}^3 e_i^{\,a_i z_4(g)},
                    \qquad 0\leq a_i<5.                       \tag{F.22.4}
\]
They have the same rationalization. The integral class-three
set injects into its rationalization, so every fibre of
rationalization on \(\mathcal S_4\), over its image, has
exactly \(125\) elements.

The same statements hold for the metabelian lower-central
quotients through class four. For their induced relative
Fermat kernels
\[
 K_{n,r}=\ker\!\left(N_n\longrightarrow
                          (\mathbb Z/5^r(1))^2\right),
                  \qquad n=3,4,\quad r\geq1,
\]
the map \(H^1(G,K_{n,r})\to\mathcal S_n\) is injective and its
image consists exactly of the classes whose six abelian
coordinates are divisible by \(5^r\). Consequently the classes
descending to every such Fermat level form the whole fibre
over the zero abelian class: a \(\mathbb Z_5^2\)-torsor in class
three, and \(125\) classes over each point of that torsor in
class four.

There is also an exact finite-quotient count. For \(n\leq4\),
reduce the integral nilpotent Lie coordinates modulo \(5^r\)
and denote the resulting characteristic quotient by \(N_{n,r}\).
Then
\[
 \begin{split}
 |H^1(G,N_{1,r})|=|H^1(G,N_{2,r})|&=5^{6r},\\
 |H^1(G,N_{3,r})|&=5^{8r},\\
 |H^1(G,N_{4,r})|&=5^{8r+3}.
 \end{split}                                                   \tag{F.22.5}
\]
The order-five class-four parameters survive every reduction
map; they are not temporary finite-quotient classes.
\end{statement}
\begin{proof}
\emph{Integral central cohomology.}
We first compute the actual integral groups needed in the
obstruction sequence:
\[
\begin{array}{c|ccc}
 j&2&3&4\\ \hline
 H^1(G,\mathbb Z_5(j))&0&\mathbb Z_5&\mathbb Z/5\\
 H^2(G,\mathbb Z_5(j))&0&0&0 .
\end{array}                                                    \tag{F.22.6}
\]
Here is a direct calculation that also checks the torsion.
Let \(L=\mathbb Q(\zeta_5)\), \(\Delta=\operatorname{Gal}(L/\mathbb Q)\),
and let \(\omega:\Delta\to\mathbb F_5^\times\) be the cyclotomic
character. The discriminant of \(L\) is \(125\). Minkowski's
bound for its ideal classes is
\[
       \frac{4!}{4^4}\left(\frac4\pi\right)^2\sqrt{125}<2.
\]
Thus \(L\) has class number one; see
\citet[Theorem~4.3]{FA_Milne2020ANT} for the precise bound.
The primes \(2\) and \(3\) are inert, since both have order
four modulo five, and \(5\) has one ramified prime.
There are therefore exactly three primes in \(S_L\), each
fixed by \(\Delta\). Put \(R_L=\mathcal O_L[1/30]\).

The Kummer sequence and the trivial ideal class group give
\[
 H^1(G_{L,S_L},\mu_5)=R_L^\times/(R_L^\times)^5,\qquad
 H^2(G_{L,S_L},\mu_5)=\operatorname{Br}(R_L)[5].
                                                               \tag{F.22.7}
\]
The comparison of these Galois groups with the corresponding
étale cohomology groups is
\citet[II, Proposition~2.9]{FA_Milne2006ADT}.
The global Brauer invariant sequence identifies the second
group with
\[
     \ker\!\left(\mathbb F_5^3\xrightarrow{\ \sum\ }\mathbb F_5\right).
                                                               \tag{F.22.8}
\]
The \(\Delta\)-action on this two-dimensional group is trivial:
each of the three primes is fixed, and the local Brauer
invariant is invariant under automorphisms of the local field.
The global sequence is recalled in
\citet[VI, \S8.1.1]{FA_Weibel2013}.

The first group in (F.22.7), as an \(\mathbb F_5[\Delta]\)-module,
is
\[
                  \mathbf1^{\oplus3}\oplus\omega\oplus\omega^2.
                                                               \tag{F.22.9}
\]
The three trivial characters are the valuations at the
three primes of \(S_L\); the valuation map is onto because
the class group is trivial. Roots of unity contribute
\(\omega\). The ordinary unit group has free rank one,
by the unit theorem. On that free part a generator
\(\delta(\zeta_5)=\zeta_5^2\) acts by \(-1\):
for the nontorsion real unit
\(\varepsilon=(1+\sqrt5)/2\), one has
\(\delta(\varepsilon)=-\varepsilon^{-1}\).
This determines the action on the entire rank-one free
lattice. Thus its character modulo five is \(\omega^2\).
The needed unit and \(S\)-unit ranks are
\citet[Theorems~5.1 and~5.11]{FA_Milne2020ANT}.
All the module extensions split for this calculation,
since \(4=|\Delta|\) is invertible in \(\mathbb F_5\).

Restriction and Hochschild--Serre now give
\[
 H^i(G,\mathbb F_5(j))
  =\left(H^i(G_{L,S_L},\mu_5)\otimes\omega^{j-1}\right)^\Delta.
                                                               \tag{F.22.10}
\]
Consequently
\[
 \begin{array}{c|rrrr}
 j&1&2&3&4\\ \hline
 \dim_{\mathbb F_5}H^1(G,\mathbb F_5(j))&3&0&1&1\\
 \dim_{\mathbb F_5}H^2(G,\mathbb F_5(j))&2&0&0&0 .
 \end{array}                                                    \tag{F.22.11}
\]
For \(j=2,3,4\), the coefficient exact sequences imply
\(H^2(G,\mathbb Z/5^r(j))=0\) for every \(r\).
All finite-coefficient cohomology groups are finite because
\(S\) is finite; this is
\citet[I, Corollary~4.15]{FA_Milne2006ADT}.
Their inverse systems have zero derived limit, and hence
\(H^2(G,\mathbb Z_5(j))=0\). Moreover
\[
 H^1(G,\mathbb Z_5(j))/5
                =H^1(G,\mathbb F_5(j)).                       \tag{F.22.12}
\]
These \(H^1\) groups are compact inverse limits of finite
groups, so (F.22.12) and the compact form of Nakayama's
lemma show that they are finitely generated over \(\mathbb Z_5\).

Their torsion is exactly
\[
          H^0(G,\mathbb Q_5/\mathbb Z_5(j)).
                                                               \tag{F.22.13}
\]
Indeed \(H^0(G,\mathbb Q_5(j))=0\), and the connecting
homomorphism for
\(0\to\mathbb Z_5(j)\to\mathbb Q_5(j)\to
\mathbb Q_5/\mathbb Z_5(j)\to0\) identifies (F.22.13)
with the kernel of rationalization on \(H^1\).
The cyclotomic image of \(G\) is \(\mathbb Z_5^\times\).
For \(j=2,3\) the group in (F.22.13) is zero; for \(j=4\)
it has order five, because all \(u^4-1\) are divisible
by five and \(v_5(2^4-1)=1\).
Combining this fact with (F.22.12) proves (F.22.6).
This calculation detects the integral torsion directly;
rational ranks alone would not determine it.

For comparison, the same vanishing is consistent with
the integral étale \(K\)-theory calculation:
\citet[VI, proof of Theorem~8.2 and Table~10.1.1]{FA_Weibel2013}
give the relevant comparison and
\(K_4(\mathbb Z)=K_6(\mathbb Z)=0\),
\(K_5(\mathbb Z)=\mathbb Z\),
\(K_7(\mathbb Z)=\mathbb Z/240\).
Localization introduces, for \(K_4\), finite residue groups
of orders \(q^2-1\), and for \(K_6\), groups of orders
\(q^3-1\), for \(q=2,3,5\). None of these orders is
divisible by five, by
\citet[IV, Corollary~1.13]{FA_Weibel2013}.
The direct argument above makes no identification of an
integral group merely from its rational dimension.

\emph{The actual group action and its fixed points.}
The tangential arithmetic action has the form
\[
 g(x)=x^{\chi_5(g)},\qquad
             g(y)=f_g^{-1}y^{\chi_5(g)}f_g,\quad f_g\in P'.
                                                               \tag{F.22.14}
\]
See \citet[formula~(9)]{FA_Wickelgren2012}.
The integral marked curve and the unit tangent have good
reduction away from five for their pro-five fundamental
group. Prime-to-residue-characteristic specialization,
as in \citet[Exposé~XIII, Corollary~2.9 and \S2.10]{FA_SGA1},
therefore makes (F.22.14) an action of the displayed \(G\).
At class two the conjugation disappears, as in F.20.
At higher classes we retain it.

The graded lower-central layers nevertheless have the
simple exact action
\[
 \gamma_jP/\gamma_{j+1}P
       =\mathbb Z_5(j)^{\ell_j(2)},\qquad
          \ell_1(2),\ell_2(2),\ell_3(2),\ell_4(2)=2,1,2,3 .
                                                               \tag{F.22.15}
\]
The free Lie algebra is generated by the abelian layer,
and (F.22.14) acts there by the scalar \(\chi_5\), which
proves the stated action in every degree. Twisting by any
section cocycle changes the action only by inner
automorphisms, which act trivially on these graded layers.
If an element in any twisted \(N_n\) were fixed, its first
nonzero graded component would be fixed by \(\chi_5^j\).
No nonzero vector in \(\mathbb Z_5(j)\) has this property:
the image of \(\chi_5^j\) is nontrivial and the lattice is
torsion free. Therefore
\[
                  H^0(G,N_n^c)=1
                  \quad\hbox{for every cocycle \(c\)}.          \tag{F.22.16}
\]
The same argument works over \(\mathbb Q_5\).

For the central extension \(N_n\to N_{n-1}\), each section
class has an obstruction in
\(H^2(G,\mathbb Z_5(n)^{\ell_n(2)})\). A continuous lift
as a function is available from Hall coordinates, so
this is the usual continuous central obstruction.
It vanishes for \(n=2,3,4\) by (F.22.6).
The lift classes form a torsor under the corresponding
\(H^1\), modulo the action of the twisted fixed group
of \(N_{n-1}\). That fixed group is trivial by (F.22.16).
Together with F.20 this proves (F.22.1)--(F.22.2).

\emph{The explicit torsion and rationalization.}
Equation (F.22.3) is integral by Fermat's congruence
and is a cocycle because
\[
 z_4(gh)=z_4(g)+\chi_5(g)^4z_4(h).
\]
Its multiple by five is the coboundary of \(1\).
It is not a coboundary: evaluation on an element with
\(\chi_5(g)=2\) would require \(3=15a\) with
\(a\in\mathbb Z_5\). Thus it generates the order-five
group in (F.22.6), proving (F.22.4).
One possible basis is
\[
 [x,[x,[x,y]]],\quad [y,[x,[x,y]]],\quad
 [y,[y,[x,y]]]\pmod{\gamma_5P}.
\]
After rationalization, (F.22.3) is the coboundary of
\(1/5\), so all \(125\) lifts become conjugate.
Conversely rationalization is injective on the
class-one set and on \(H^1(G,\mathbb Z_5(3))\); the
class-two fibres are singletons. Apply the central
cohomology sequence inductively, using the rational
version of (F.22.16), to see that it is injective on
\(\mathcal S_3\). There can therefore be no additional
identifications between the displayed class-four fibres.

\emph{Metabelian quotients and relative Fermat descent.}
For a rank-two free group,
\([P',P']\subset\gamma_5P\): modulo \(\gamma_5P\) the
only potential degree-four bracket is the bracket of
the rank-one degree-two layer with itself, which is
zero. Equivalently, write elements of \(P'\) modulo
\(\gamma_3P\) as powers of the same commutator and use
\([\gamma_2P,\gamma_3P]\subset\gamma_5P\).
Thus imposing metabelianity changes none of the
quotients under discussion.

For \(D_r=(\mathbb Z/5^r(1))^2\), one has
\(H^0(G,D_r)=0\), since \(2-1\) is a unit modulo five.
The same is true after every twist arising from
\(K_{n,r}\), which acts trivially on \(D_r\).
The nonabelian cohomology sequence for
\(1\to K_{n,r}\to N_n\to D_r\to1\) consequently makes
the first \(H^1\) map injective, with image the inverse
image of the zero class in \(H^1(G,D_r)\).
Kummer theory identifies this last map with reduction
of the six abelian coordinates modulo \(5^r\).
This proves all the relative Fermat assertions, including
the description of the intersection over every \(r\).
The quotients are induced from \(P\); no assertion about
the absolute metabelian completion of a Fermat cover
is substituted for this calculation.

\emph{Finite quotients and exact checks.}
For class at most four, all the denominators in the
Baker--Campbell--Hausdorff and inverse Hall-coordinate
formulas are units in \(\mathbb Z_5\). The integral
nilpotent Lie lattice therefore defines \(N_n\), and
its reduction modulo \(5^r\) defines \(N_{n,r}\).
The full arithmetic action preserves this lattice
and its characteristic sublattices \(5^r\).
From (F.22.6) and the coefficient sequence one obtains
\[
 \begin{split}
 H^1(G,\mathbb Z/5^r(2))&=0,\\
 H^1(G,\mathbb Z/5^r(3))&=\mathbb Z/5^r,\\
 H^1(G,\mathbb Z/5^r(4))&=\mathbb Z/5,
 \end{split}\qquad
 H^2(G,\mathbb Z/5^r(j))=0\quad(j=2,3,4).
                                                               \tag{F.22.17}
\]
The last \(H^1\) is the reduction of the integral group
of order five, so its transition maps are isomorphisms.
The twisted fixed groups in the bases \(N_{1,r}\),
\(N_{2,r}\), and \(N_{3,r}\) are all trivial: for their
degrees \(j=1,2,3\), the numbers \(2^j-1\) are units
modulo five. The finite central fibres thus have
orders \(1\), \(5^{2r}\), and \(5^3\), respectively.
This proves (F.22.5).

The accompanying exact program verifies the discriminant
and the strict Minkowski bound, the two inertia orders,
the character multiplicities in (F.22.11), and every
residue order in the \(K\)-theory check. It also evaluates
(F.22.3) on the cyclic groups
\((\mathbb Z/5^{r+1})^\times\), for \(1\leq r\leq6\),
checks its cocycle equation on a generator and the
wraparound relation, and verifies that its class
generates the quotient by all coboundaries.
The finite section counts are recorded as consequences
of the proved cohomology calculation, rather than as
an enumeration of a finite presentation for \(G\).
The higher unrestricted global stages in F.12 remain
outside the scope of this finite-ramification result.
\end{proof}

\subsection{The first integral obstruction and the relative Fermat fibre}

\begin{statement}{F.23}{Proved}{Actual class-five congruences and complete relative Fermat fibres}
Retain \(G,N_n,\mathcal S_n\) from F.22. Put
\[
 C=[X,Y],\quad U=[X,C],\quad V=[Y,C],\quad
 E_1=[X,U],\quad E_2=[Y,U],\quad E_3=[Y,V]
\]
in the free Lie algebra on \(X=\log x,Y=\log y\), and use the
degree-five basis
\[
 (B_1,\ldots,B_6)=
 ([X,E_1],[Y,E_1],[Y,E_2],[Y,E_3],[C,U],[C,V]).
                                                               \tag{F.23.1}
\]
There are exact central cohomology groups
\[
 H^1(G,\mathbb Z_5(5))=\mathbb Z_5,\qquad
 H^2(G,\mathbb Z_5(5))=(\mathbb Z/5)^2,                         \tag{F.23.2}
\]
and restriction identifies the second group with the direct
sum of its local groups at \(2\) and \(3\).

For any actual class-four cocycle, define its three local
coordinates \(v_q=(v_{q,1},v_{q,2},v_{q,3})\in\mathbb F_5^3\)
at \(q=2,3\) by the explicit normalization below.
Let \(a_q,b_q\in\mathbb F_5\) be its two Kummer valuation
coordinates at \(q\). The cocycle lifts to class five if and
only if, at both primes, the following six equations hold:
\[
\begin{aligned}
 a v_1&=a^4b,&
 b v_1+a v_2&=3b-a^3b^2,&
 b v_2+a v_3&=2b+a^2b^3,\\
 b v_3&=3b-ab^4,&
 a v_2&=a^3b^2,&
 a v_3&=2a^2b^3 .
\end{aligned}                                                   \tag{F.23.3}
\]
Here \(a=a_q,b=b_q,v=v_q\), and all equations are in
\(\mathbb F_5\). Their complete solution set is
\[
\begin{array}{c|c}
(a,b)&v\\ \hline
(0,0)&\text{arbitrary in }\mathbb F_5^3\\
(a,0),\ a\ne0&(0,0,0)\\
(0,b),\ b\ne0&(3,2,3)\\
(a,a),\ a\ne0&(1,1,2)\\
ab(a-b)\ne0&\text{no solution}.
\end{array}                                                     \tag{F.23.4}
\]
Each nonempty global class-five fibre is a free
\(\mathbb Z_5^6\)-torsor. In the metabelian quotient,
the last two equations are omitted, but the solution set
(F.23.4) is unchanged; its nonempty fibres are
\(\mathbb Z_5^4\)-torsors. In particular, every class-four
section above the abelian pair
\((\kappa(2),2\kappa(2))\) fails to lift, both in the full
and in the metabelian tower.

All the lower global parameters are retained in this test.
For a fixed global class-three class, choose any class-four
lift and write its two normalized local vectors as
\(v_2^0,v_3^0\). The \(125\) class-four lifts are indexed by
\(t\in\mathbb F_5^3\), and their local vectors are exactly
\[
                       v_2^0+t,\qquad v_3^0+t.                \tag{F.23.5}
\]
Thus when both nonzero valuation pairs are among the three
permitted directions, a class-four lift survives precisely
when
\[
 v_3^0-v_2^0=R(a_3,b_3)-R(a_2,b_2),                          \tag{F.23.6}
\]
where \(R\) is the corresponding singleton vector in
(F.23.4). When exactly one valuation pair is nonzero in a
permitted direction, exactly one of the \(125\) lifts
survives. When both valuation pairs are zero, all \(125\)
survive. The difference on the left of (F.23.6) is independent
of the chosen class-four origin; its dependence on the full
global class-three cocycle is not discarded.

In particular, the entire induced relative Fermat tower at
class five is evaluated. For \(r\geq1\), let
\[
 K_{5,r}=\ker\!\left(N_5\longrightarrow
                              (\mathbb Z/5^r(1))^2\right).
\]
After choices of torsor origins,
\[
\begin{aligned}
 H^1(G,K_{5,r})&\simeq
        (5^r\mathbb Z_5)^6\times\mathbb Z_5^8
                                      \times(\mathbb Z/5)^3,\\
 H^1(G,K_{5,r}^{\rm met})&\simeq
        (5^r\mathbb Z_5)^6\times\mathbb Z_5^6
                                      \times(\mathbb Z/5)^3.
\end{aligned}                                                   \tag{F.23.7}
\]
The maps to the ambient section sets are injective and
identify transition maps with inclusion of the indicated
abelian divisibility loci. The fibre over abelian zero,
equivalently the classes descending to every relative
Fermat level, is respectively
\[
 \mathbb Z_5^8\oplus(\mathbb Z/5)^3,\qquad
 \mathbb Z_5^6\oplus(\mathbb Z/5)^3.                           \tag{F.23.8}
\]
For this last assertion an explicit integral diagonalization
of the actual arithmetic action is given in the proof.
These are induced relative quotients and retain the stated
finite ramification set.
\end{statement}
\begin{proof}
\emph{The arithmetic obstruction group and its local detection.}
First replace \(S\) by \(S_0=\{5\}\). Over
\(\mathbb Q(\zeta_5)\), the ring with just the prime over
five inverted has trivial ideal class group and trivial
Brauer group: in the global Brauer sequence there is only
one finite place and no real place. The argument of F.22
therefore gives
\[
 H^2(G_{\mathbb Q,S_0},\mathbb Z_5(j))=0\quad(j\geq1).
                                                               \tag{F.23.9}
\]
At weight five, the \(S_0\)-unit module modulo fifth powers
has characters \(\mathbf1\oplus\omega\oplus\omega^2\).
Twisting by \(\omega^4=1\) selects one dimension.
The torsion invariant
\(H^0(G_{\mathbb Q,S_0},\mathbb Q_5/\mathbb Z_5(5))\)
is zero. The same compact Nakayama argument as in F.22
gives \(H^1(G_{\mathbb Q,S_0},\mathbb Z_5(5))=\mathbb Z_5\).

The localization sequence on \(\operatorname{Spec}\mathbb Z[1/5]\)
when \(2,3\) are deleted has residue terms
\(H^{i-1}(\mathbb F_q,\mathbb Z_5(4))\).
For completeness the codimension-one purity calculation
needed here is the elementary one in
\citet[II, \S0.1(f), Corollary~1.6 and Remark~1.7(b)]{FA_Milne2006ADT}:
\(R^2i^!\mu_{5^r}=\mathbb Z/5^r\), with the other
cohomology sheaves zero. Tensoring by the unramified
fourth twist gives the displayed residue terms.
Passing to integral coefficients is exact because every
finite-coefficient group in this fixed finite set is finite.
The residue-field invariants in the integral module
\(\mathbb Z_5(4)\) are zero, and
\[
 H^1(\mathbb F_q,\mathbb Z_5(4))
           =\mathbb Z_5/(q^4-1).
\]
Also \(H^3(G_{\mathbb Q,S_0},\mathbb Z_5(5))=0\), by odd
cohomological dimension two and the same inverse-limit
argument. Thus localization gives (F.23.2), since
\(v_5(2^4-1)=v_5(3^4-1)=1\), and does not change \(H^1\).
The local inertia--residue exact sequence identifies
\(H^2(\mathbb Q_q,\mathbb Z_5(5))\) with this same residue
group. This proves that the two local restrictions detect
the entire global obstruction. No claim that the entire
group \(K_8(\mathbb Z)\) is known to vanish is needed.

\emph{The actual Frobenius action through the required degree.}
Fix tame inertia \(\tau_q\) and arithmetic Frobenius
\(\phi_q\), so that
\(\phi_q\tau_q\phi_q^{-1}=\tau_q^q\).
The geometric arithmetic action is unramified at \(q=2,3\).
Write the logarithm of its actual associator as
\[
 \log f_{\phi_q}
 =r_qC+h_q(U+V)+c_q(E_1+E_3)+d_qE_2
                                      \pmod{\text{degree }5}.
                                                               \tag{F.23.10}
\]
The two-cycle gives the indicated equal coefficients;
the class-two hexagon in E.16 gives
\(r_q=(q^2-1)/24\).
The degree-three coefficient \(h_q\) is retained. It is
the first Soulé parameter, up to its stated normalization
and possible coboundary, not a zero coefficient assumed
from a graded model.

The actual pentagon forces the following exact values:
\[
\begin{array}{c|rrr}
 q&r_q&c_q&d_q\\ \hline
 2&1/8&-1/96&-1/128\\
 3&1/3&-1/18&-1/36.
\end{array}                                                     \tag{F.23.11}
\]
Here is an explicit certificate for this step.
Use the action of \(PB_4\) on its free fibre with generators
\(u=x_{14},v=x_{24},w=x_{34}\), as in the presentation
checked in E.17. For \(A=x_{12},B=x_{23},D=x_{13}\),
the conjugations \(g^{-1}(\,\cdot\,)g\) are
\[
\begin{array}{c|ccc}
 &u&v&w\\ \hline
 A&uvu v^{-1}u^{-1}&uvu^{-1}&w\\
 B&u&vwv w^{-1}v^{-1}&vwv^{-1}\\
 D&uwu w^{-1}u^{-1}&[u,w]v[u,w]^{-1}&uwu^{-1}.
\end{array}                                                     \tag{F.23.12}
\]
The fibre generators act by their own inner conjugations.
The program verifies the nine displayed fibre relations
and the three centrality relations for \(ADB\), and
evaluates the actual group pentagon, with its five
insertion/cabling maps, in this representation.
Faithfulness is unnecessary: every genuine pentagon
must hold in this representation.

In the image of the generator \(u\), the coefficients
of the Magnus words \(u^3vw\) and \(u^3wv\), where these
letters now denote the augmentation variables, give
the two equations
\[
 3c_q-2d_q=-r_q^2,\qquad
       -c_q-d_q=r_q/12+r_q^2/2.                              \tag{F.23.13}
\]
Their coefficient determinant is \(-5\).
All remaining associative coefficients of the pentagon
and hexagon are checked by exact rational arithmetic.
Terms in \(h_q\) enter affinely and cancel; products of
a degree-three term with a degree-two term first occur
beyond the degree being checked. Solving (F.23.13) in
\(\mathbb Q_5\), then checking the resulting denominators
are five-adic units for these two \(q\), proves (F.23.11).
The nonunit determinant is kept in the certificate.

The parameter \(h_q\) can now be removed \emph{locally by
an integral change of geometric coordinates}.
Let \(\eta_q\) fix \(x\) and send
\[
 y\longmapsto k_q^{-1}yk_q,\qquad
 \log k_q=-\frac{h_q}{q^3-1}(U+V).
                                                               \tag{F.23.14}
\]
These are integral elements through class five:
\(q^3-1\) is a five-adic unit, and the logarithm of the
subgroup \(\gamma_3N_5\) has only five-integral
denominators, as checked explicitly below.
Conjugating the action by \(\eta_q\) in the form
\(\eta_q^{-1}\rho(\phi_q)\eta_q\) makes \(h_q=0\) in
(F.23.10) and preserves \(r_q,c_q,d_q\).
The program checks this exact equality through class
five; degree considerations make it an affine identity
in \(h_q\). The section cocycle is transported by
\(\eta_q^{-1}\) at the same time. Thus its arithmetic
information has been retained in its coordinates.

\emph{Definition of the local coordinates and the twelve equations.}
After (F.23.14), conjugate a local class-four cocycle
so that its value on \(\phi_q\) has no components of
degrees one, two or three. This is possible successively
and uniquely in those degrees because \(q^j-1\) is a
five-adic unit for \(j=1,2,3\).
Its remaining value is in the central fourth layer;
write its logarithm there as \(W_q=\sum_i w_{q,i}E_i\).
Changing the fourth component of the conjugating element
changes \(W_q\) by \((q^4-1)L_4\). Therefore
\[
 v_q=\left(\frac{q^4-1}{5}\right)^{-1}
                    (w_{q,1},w_{q,2},w_{q,3})\pmod5          \tag{F.23.15}
\]
is well defined for the local section class and the
declared geometric coordinates.

Its inertia value \(I_q\in N_4\) has abelianization
\(aX+bY\), with the integral Kummer valuations \(a,b\).
Because the Frobenius value is central in \(N_4\),
the tame cocycle relation is
\[
                   \rho(\phi_q)(I_q)=I_q^q\quad\text{in }N_4.
                                                               \tag{F.23.16}
\]
For any given \(a,b\in\mathbb Z_5\), this determines the
inertia value uniquely: in degrees \(j=2,3,4\), correcting
the \(j\)th Hall coordinate multiplies its defect by
\(q^j-q\), which is a five-adic unit.
The program implements exactly this recursive solution,
starting with \(x^ay^b\), and retains all the resulting
Hall coordinates.

Lift the resulting word to \(N_5\) with its six fifth
Hall coordinates set to zero. The fifth-degree defect
of (F.23.16), divided by \(q(q^4-1)/5\), is
\[
 P(a,b)=
 (a^4b,\ 3b-a^3b^2,\ 2b+a^2b^3,\ 3b-ab^4,\
                                  a^3b^2,\ 2a^2b^3)
                      \quad\text{in }\mathbb F_5^6.           \tag{F.23.17}
\]
This is an actual Magnus calculation, with the coefficients
(F.23.11) inserted. It uses group products and power series
with integral coefficients, rather than applying a
class-less-than-five BCH correspondence in class five.
The first four Hall-coordinate corrections have denominators
prime to five. Dependence on the abelian exponents modulo
five is sufficient: modulo five, replacing \(x^a\) by
\(x^{a+5}\) changes the Magnus series first by the central
degree-five term \(X^5\); its contributions to the two
sides of (F.23.16) are \(q^5X^5\) and \(qX^5\), which
agree modulo five. The same holds for \(Y^5\); fifth
powers of commutators first occur in degree at least ten.
Thus all \(25\) residue pairs exhaust this calculation.
The exact program verifies (F.23.17) on each pair for
both Frobenius values.

Restoring the fourth-layer Frobenius value contributes
the commutator \(q[W_q,aX+bY]\) to the defect. In the basis
(F.23.1), its negative, after the same normalization, is
\[
 (av_1,\ bv_1+av_2,\ bv_2+av_3,\ bv_3,\ av_2,\ av_3).
                                                               \tag{F.23.18}
\]
These formulas follow directly from Jacobi:
\([X,E_2]=B_2+B_5\) and \([X,E_3]=B_3+B_6\).
The local obstruction is therefore exactly
(F.23.3). Changing the fifth-layer inertia lift alters
the integral defect by \((q^5-q)L_5\); changing the
central Frobenius lift has no effect. Since
\(v_5(q^5-q)=1\), vanishing of the displayed residue
is necessary and sufficient for a local lift.
Wild inertia and the other prime-primary tame factors
have trivial maps into the pro-five target. Conversely
the tame relation suffices, since the Frobenius and
inertia values then define a continuous cocycle on the
tame local quotient.

The solution table (F.23.4) can also be checked without
enumeration. On the two coordinate axes the equations
immediately give the displayed vectors. If \(ab\ne0\),
the first, fifth and sixth equations give
\(v_1=a^3b,v_2=a^2b^2,v_3=2ab^3\).
The second then forces \(a^3b=1\), hence \(a=b\) in
\(\mathbb F_5^\times\), and the remaining equations hold.
For the metabelian quotient, \(B_5,B_6\) vanish.
Solving just the first four equations when \(ab\ne0\)
reduces, with \(t=b/a\), to
\((t+1)^3=3\). Cubing permutes \(\mathbb F_5^\times\);
the unique solution is \(t=1\).
The exact enumeration gives \(125+12=137\) permitted
parameter pairs out of \(25\cdot125=3125\) at each prime,
for both the full and metabelian calculation.

By the localization isomorphism already proved, these
two local tests are the complete global obstruction test.
The fixed-group argument of F.22 works integrally in
every finite lower-central class, including the twisted
class-four base. Thus nonempty fibres are free torsors
under \(H^1(G,L_5)\), of ranks six and four as claimed.
The example \((\kappa(2),2\kappa(2))\) has residue pair
\((1,2)\) at \(2\), which fails (F.23.4), independently
of all of its higher choices.

\emph{The common shift and the retained global parameters.}
For a fixed global class-three class the \(125\) lifts
are exactly the twists by
\(z_4(g)\sum_i t_iE_i\) of F.22. At \(\phi_q\) this
adds \((q^4-1)t/5\). The change (F.23.14) acts trivially
on the fourth layer, and normalization in lower degrees
is unchanged. Hence (F.23.15) changes by the same vector
\(t\) at both primes. This proves (F.23.5)--(F.23.6).
In particular the complete global class-three parameters
are retained inside \(v_3^0-v_2^0\); independently
normalizing both local vectors to zero would change
the global problem and is not done.

When all six abelian coordinates are divisible by
\(5^r\), \(r\geq1\), both residue pairs are zero.
Every class-four section in this relative Fermat image
therefore lifts. The injection and image argument for
the Fermat kernels in F.22 applies unchanged to \(N_5\).
It yields (F.23.7), including the transition maps.

\emph{An integral diagonalization of the entire abelian-zero fibre.}
The subgroup \(M=\gamma_3N_5\) is abelian, since brackets
of two third-layer elements have degree at least six.
Its logarithm is the integral lattice
\[
                          L_3\oplus L_4\oplus L_5.
                                                               \tag{F.23.19}
\]
This does not assert an integral BCH correspondence for
all of \(N_5\). To verify (F.23.19), take the two third,
three fourth and six fifth Hall commutators specified
above. Their logarithms, in increasing weights, form a
unipotent triangular change of basis whose entries have
denominators prime to five. Products in \(M\) add their
logarithms. The exact \(11\)-by-\(11\) change-of-basis
matrix is retained in the certificate. Thus this
particular abelian subgroup has the stated integral
logarithm lattice even in class five.

For arbitrary \(g\in G\), put \(\lambda=\chi_5(g)\).
Only the quadratic term
\((\lambda^2-1)C/24\) of \(\log f_g\) acts nontrivially
on \(L_3\) through degree five. Define the degree-two
derivation
\[
 D_2(X)=0,\qquad D_2(Y)=[Y,C].
\]
Then for \(v\in L_3\),
\[
 \rho(g)v=\lambda^3
       \left(v+\frac{\lambda^2-1}{24}D_2(v)\right),
 \qquad
 \rho(g)|_{L_4}=\lambda^4,\quad
 \rho(g)|_{L_5}=\lambda^5.                                   \tag{F.23.20}
\]
This follows by expanding the actual action
\(\rho(g)Y=\exp(-\operatorname{ad}\log f_g)(\lambda Y)\).
In particular
\[
 D_2(U)=B_2+B_5,\qquad D_2(V)=B_3-B_6.
\]
The integral corrected basis
\[
        \widetilde U=U-(B_2+B_5)/24,\qquad
        \widetilde V=V-(B_3-B_6)/24                           \tag{F.23.21}
\]
therefore transforms by \(\lambda^3\). This proves the
actual \(G\)-module isomorphism
\[
 M\simeq \mathbb Z_5(3)^2\oplus
          \mathbb Z_5(4)^3\oplus\mathbb Z_5(5)^6 .
                                                               \tag{F.23.22}
\]
The unknown higher Soulé coefficients have degree too
large to affect this calculation; they have not been
assigned artificial values in the action.
In the metabelian quotient discard \(B_5,B_6\), giving
four last summands.

Since \(H^1(G,N_2)\to H^1(G,N_1)\) is a bijection and
every twisted \(N_2\) has trivial fixed group,
\(H^1(G,M)\to H^1(G,N_5)\) is injective with image
exactly the abelian-zero fibre. Taking cohomology in
(F.23.22), using F.22 and (F.23.2), proves (F.23.8).
For chosen integral cocycles \(\eta_3,\eta_5\) generating
the two cyclic free cohomology groups, its elements
are explicitly represented by
\[
 \exp\!\left(
 \eta_3(g)(s_1\widetilde U+s_2\widetilde V)
  +z_4(g)\sum_{i=1}^3t_iE_i
  +\eta_5(g)\sum_{j=1}^6r_jB_j\right),
                                                               \tag{F.23.23}
\]
where \(s_i,r_j\in\mathbb Z_5\) and \(0\leq t_i<5\).
The exponential in this formula is the integral
isomorphism for the abelian group \(M\).

The finite certificate supplies all polynomial identities,
the nonzero pentagon minor, local coordinate changes,
all \(6250\) local parameter tests, and the integral
matrix in (F.23.19). Formula (F.23.6) remains expressed
in the actual localized global cocycles: a separate
closed formula for that frame difference in arbitrary
chosen global torsor coordinates has not been asserted.
The explicit no-lift example and the entire relative
Fermat calculation do not require such a formula.
\end{proof}

\begin{statement}{F.24}{Proved}{The unit cusps first separate at integral class three}
For the actual pro-five action, the three unit tangential
sections at \(0+1\), \(1-1\), and the image of \(0+1\)
under \(z\mapsto1/z\) have the same class-one and
class-two classes, but are pairwise distinct in class
three. This holds over the finite ramification set
\(S=\{2,3,5\}\), and also after inflation to \(G_{\mathbb Q}\).
All three classes descend to every induced relative
Fermat level, as described in F.22--F.23.

More precisely, use the standard unit-interval path
and the logarithmic path convention of
\citet[\S2.1, (2.4), Proposition~2.2 and Lemma~2.4]{FA_NakamuraSakugawaWojtkowiak2017}.
For \(\sigma\in\ker\chi_5\), its path cocycle satisfies
\[
       (\log f_\sigma)_3
            =\frac{\kappa_3(\sigma)}{48}
                         [X+Y,[X,Y]],                       \tag{F.24.1}
\]
where \(\kappa_3\) is the raw cyclotomic Soulé character
defined by the Kummer products in that source.
Its reduction modulo five is nonzero; in fact
\(\kappa_3\) has image \(\mathbb Z_5\) on the
cyclotomic kernel. The path convention is part of
(F.24.1); reversing a transition path changes the
sign of its first nonzero term and does not affect
the separation assertion.
\end{statement}
\begin{proof}
The class-two collision was proved in F.20. We now
verify an integral nonzero third-layer term, including
its normalization and nonvanishing.

Put \(L=\mathbb Q(\zeta_5)\), and
\(\varepsilon=(1+\sqrt5)/2\).
The first Kummer product for \(\kappa_3\) is
\[
               \epsilon_{3,1}
                   =\prod_{a=1}^4(1-\zeta_5^a)^{a^2}.
\]
Modulo fifth powers in \(L^\times\), pair the factors
with indices \(a\) and \(5-a\) to obtain
\[
\begin{split}
 \epsilon_{3,1}
 &\equiv
 \bigl((1-\zeta_5)(1-\zeta_5^{-1})\bigr)
 \bigl((1-\zeta_5^2)(1-\zeta_5^{-2})\bigr)^4\\
 &=\left(\frac{\sqrt5}{\varepsilon}\right)
                   (\sqrt5\,\varepsilon)^4
   =(\sqrt5)^5\varepsilon^3
                 \pmod{L^{\times5}}.                        \tag{F.24.2}
\end{split}
\]
The unit \(\varepsilon\) is not a fifth power in \(L\).
Indeed it is the fundamental positive unit of
\(\mathbb Q(\sqrt5)\). One can verify the latter
elementary claim directly: a positive unit strictly
between \(1\) and \(\varepsilon\) would have integral
trace strictly between \(0\) and \(1\) if its norm is
\(-1\), or strictly between \(2\) and \(3\) if its norm
is \(1\). The rank-one unit theorem then makes
\(\varepsilon\) a generator. If \(u^5=\varepsilon\)
in \(L\), taking the norm to its real subfield would
give \(N(u)^5=\varepsilon^2\), which is impossible
in that unit group.

Since \(L\) contains \(\mu_5\), Kummer theory now gives
\([L(\varepsilon^{1/5}):L]=5\).
Furthermore \(L(\varepsilon^{1/5})\) is not contained
in \(\mathbb Q(\mu_{5^\infty})\). It is Galois over
\(\mathbb Q\): the nontrivial real-subfield
automorphism sends \(\varepsilon\) to
\(-\varepsilon^{-1}\), and \(-1\) has a fifth root
in \(\mathbb Q\). Complex conjugation fixes a real
fifth root of \(\varepsilon\), but inverts the
nontrivial Kummer group of order five. Thus this
Galois extension is nonabelian, whereas every finite
subextension of the cyclotomic extension is abelian.
Consequently (F.24.2) gives a nonzero Kummer character
even after restriction to the full cyclotomic kernel.
The defining congruences for the Soulé character,
at \(n=1\), identify it with the reduction of
\(\kappa_3\) modulo five. This proves the asserted
nonzero reduction. Its compact image is a closed
subgroup of \(\mathbb Z_5\) containing a unit, hence
is all of \(\mathbb Z_5\).

For the coefficient in (F.24.1), the cited source
defines the path cochain by
\(\gamma=f_\sigma*\sigma(\gamma)\) and writes
\[
 \log f_\sigma\equiv-\kappa_z(\sigma)X
       -\sum_{m\geq1}\ell i_m(z)(\sigma)
                         \operatorname{ad}_X^{m-1}Y
\]
in the quotient retaining words with at most one \(Y\).
Although its main application is local, the paragraph
following Proposition~2.2 explicitly gives the same
Kummer measure over number fields, and (2.11)
specializes it to \(F=\mathbb Q,z=1\). Thus the
identity used here is an identity on the global
cyclotomic kernel.
At the unit tangent \(z=1\), Proposition~2.2 and
Lemma~2.4 give, on the cyclotomic kernel,
\[
 \ell i_3(1)=\frac{\widetilde\kappa_3}{2},
 \qquad
 \widetilde\kappa_3=\frac{\kappa_3}{1-5^2}.
                                                               \tag{F.24.3}
\]
The same paper's (2.11) explicitly identifies this
restricted Kummer character with the usual Soulé
character when \(z=1\). Equations (F.24.3) make the
coefficient of \([X,[X,Y]]\) in \(\log f_\sigma\)
equal to \(\kappa_3/48\).
The two-cycle exchanges \(X\) and \(Y\) and forces
the same coefficient on \([Y,[X,Y]]\), proving
(F.24.1). This direct path computation also agrees
with the independent Tate substitution in
\citet[Theorems~5.2.4--5.2.6]{Nakamura2013Eisenstein}.

For \(\sigma\in\ker\chi_5\), the quadratic associator
coefficient vanishes by E.16. Thus \(f_\sigma\in\gamma_3P\),
and the actual arithmetic action on \(N_3\) is trivial:
conjugating a degree-one generator by \(f_\sigma\)
first changes degree four. A coboundary on \(G\),
when restricted to this subgroup, must therefore
be the identity. The transition cocycle between the
two unit tangents has the nonzero restriction
(F.24.1), so their class-three section classes
cannot be equal. The same argument works over
\(G_{\mathbb Q}\), since the Kummer extension used
above is unramified away from five.

Finally the rational \(S_3\)-symmetry of the tripod
acts on arithmetic section classes and permutes the
three cusps. On the chosen unit tangents it changes
a tangent parameter at most by a sign.
Changing a cusp tangent by \(-1\) does not change its
pro-five section class: its Kummer class is zero,
as the compatible fifth-power roots of \(-1\)
can all be chosen to be \(-1\) in \(\mathbb Q\).
The three cusp packets are consequently permuted
transitively on ordered pairs. Separation of the
first two proves separation of every pair.
Their abelian Kummer coordinates are all zero, so
the relative Fermat descent assertions of F.22--F.23
apply to each of them.

The certificate for F.23 also verifies the exact
cyclotomic identity (F.24.2), retaining a fifth-power
factor rather than only a numerical congruence.
For comparison, it records
\(\epsilon_{5,1}\equiv5\pmod{L^{\times5}}\).
Here \(5\) is not a fifth power in \(L\), since its
valuation at the prime over five is \(4\).
The same real Kummer and cyclotomic argument shows
that the raw \(\kappa_5\) is primitive. Reindexing
the defining Kummer product by multiplication by
\(\chi_5(g)\) gives
\(\kappa_m(g\sigma g^{-1})=\chi_5(g)^m\kappa_m(\sigma)\)
for \(m=3,5\) and \(\sigma\in\ker\chi_5\).
These characters are unramified away from five,
since all the defining factors are cyclotomic
five-units and their Kummer extensions are obtained
by adjoining fifth-power roots. The characters
\(\kappa_3,\kappa_5\) extend uniquely as cohomology
classes in \(H^1(G_{\mathbb Q,\{5\}},\mathbb Z_5(3))\)
and \(H^1(G_{\mathbb Q,\{5\}},\mathbb Z_5(5))\):
in inflation--restriction the cyclotomic quotient
has zero \(H^1\) and \(H^2\) with either coefficient,
because its order-four subgroup acts nontrivially.
They can therefore be used for the primitive
generators \(\eta_3,\eta_5\) in (F.23.23), with
the indicated explicit normalization.
\end{proof}


\subsection{Exact lifting statements and the remaining section problems}

\begin{statement}{F.10}{Proved}{The obstruction at each successive central extension}
Let $1\to A\to U\to V\to1$ be a central extension of
topological $G_k$-groups, with $A$ abelian and with
continuous local lifts for the cocycles under consideration.
For a $V$-cocycle $c$, choose a continuous lift $\widetilde c$.
Its defect
\[
 o_c(\sigma,\tau)=
 \widetilde c_\sigma\,\sigma(\widetilde c_\tau)
                    \widetilde c_{\sigma\tau}^{-1}\in A
\]
defines a class in $H^2_{\rm cts}(G_k,A)$ independent
of the lift. The class vanishes exactly when $c$ lifts.
For a fixed representative $c$, its lifts modulo
$A$-conjugacy form a torsor under $H^1_{\rm cts}(G_k,A)$.
Passing to $U$-conjugacy further divides by the stabilizer
of $c$ in $V$. Local conditions impose the corresponding
local obstruction and localization conditions.
\end{statement}
\begin{proof}
Associativity and the cocycle equation in $V$ give the
usual degree-two cocycle identity for $o_c$. Centrality
ensures that no additional conjugation on $A$ occurs.
Replacing $\widetilde c$ by $a_\sigma\widetilde c_\sigma$
changes its defect by the coboundary of the $A$-cochain
$a$. Vanishing therefore permits correction to an
actual cocycle, and the quotient of two corrected
lifts is an $A$-cocycle. Quotienting by conjugation gives
the stated torsor and stabilizer action.
For unipotent groups over $\Q_p$, exponential coordinates
supply the required continuous lifts. For profinite
extensions this local-lift hypothesis must be checked,
or the assertion formulated using continuous extensions
and their nonabelian cohomology instead.
Applying the construction at every lower-central stage
gives the exact successive test proposed at scans173
and178. Merely naming its $H^1$ and $H^2$ groups does
not evaluate them.
\end{proof}

\begin{statement}{F.11}{Proved}{The condition for restricting a section to a cover}
Suppose $1\to N\to P\to H\to1$ is an exact geometric
finite-cover sequence, stable under the arithmetic
action after a specified finite extension of the base
field. If the chosen approximation preserves this
exact sequence, a section of its arithmetic extension
restricts to the cover exactly when its induced section
of $H\rtimes G_k$ is the section specifying that cover,
up to the stated $H$-conjugacy. Restriction after an
open subgroup of $G_k$ is possible whenever the
corresponding finite cocycle becomes trivial there.
\end{statement}
\begin{proof}
The arithmetic group of a fixed connected cover is
the inverse image of its finite section in the
quotient extension. A section lands in that inverse
image precisely under the displayed condition.
Changing the finite section by conjugation changes
the chosen lift of the base point and yields exactly
the stated ambiguity. A continuous finite cocycle
and the finite action factor through a finite
quotient of $G_k$; their common open kernel kills
that cocycle. Exactness then identifies the geometric
kernel with the approximation of $N$, so the
restricted section is a section for the cover.
If that comparison is not an isomorphism, the
same conclusion for the cover's absolute approximation
does not follow.

For the Fermat tower B.5 applies to the virtually
abelian envelope system, and pro-$\ell$ completion
applies exactly to its $\ell$-power levels. Abelian,
metabelian, bounded nilpotent, and absolute
prounipotent completion have the failures proved
in B.7--B.9. In particular a Selmer point on the
base cannot simply be treated as a section of the
absolute completion of every Fermat cover.
\end{proof}

\begin{statement}{F.12}{Open}{The unrestricted Fermat section sets}
The entire unrestricted global section sets for the exact
Fermat approximation systems and the higher integral
metabelian quotients have not been evaluated.
The abelian groups in F.4, the rational mixed-Tate
Selmer towers in F.2, the full real pro-$2$ metabelian
section set in F.19, and the global integral pro-$5$
class-two section sets in F.20--F.21 are computed.
The latter include unrestricted ramification and its
derived-limit obstruction. F.22 computes the complete
integral class-three and class-four sets over
\(G_{\mathbb Q,\{2,3,5,\infty\}}\), including their
order-five torsion. These results also evaluate
their specified induced relative Fermat kernels.
F.23 gives the exact class-five obstruction at this
finite ramification set, a concrete class-four section
with no lift, and the complete relative Fermat
class-five sets. F.24 proves that the three specified
unit cusp sections first separate in integral class
three. These calculations
do not identify those kernels with the absolute
metabelian completions of the covers.
\end{statement}
Here is the precise residual problem and the attack made.
After fixing an arithmetic action and a reference section,
one must compute
\[
 H^1_{\rm cts}(G_k,P_{\rm met}),\qquad
 H^1_{\rm cts}(G_k,P^{(\ell)}),\qquad
 H^1_{\rm cts}(G_k,\mathcal V_{\rm ar}),
\]
with the field, packet, and finite-cover descent conditions
of F.11; $\mathcal V_{\rm ar}$ denotes the chosen arithmetic
realization of the virtually abelian envelope system.
The bare pro-discrete group $V(F)$ by itself does not
carry a specified action of all of $G_k$: cyclotomic
scalars already require profinite or adelic coefficient
groups. Its arithmetic realization must therefore be
part of the data, as in B.1, before this expression is
defined.

For the metabelian quotient, the first extension has
kernel the completed Alexander module, acted on
nontrivially through abelianization. The calculation of
that module in B.23 is explicit. F.19 evaluates its
three allowed real twists. The full global
Alexander-module $H^1$ and obstruction $H^2$ over
$\mathbb Q$ for every abelian cocycle remain unevaluated.
F.20 computes the first global class-two groups with
finite ramification; F.21 removes that restriction and
gives the exact additional completion condition.
F.22 then evaluates the next two integral layers for
\(S=\{2,3,5\}\), with free \(\mathbb Z_5^2\)-torsor
fibres in class three and exactly \(125\) lift classes
in class four. F.23 evaluates the next obstruction
by twelve explicit local congruences and retains the
common global shift in the fourth-layer parameters.
It evaluates all positive-level induced relative
Fermat sets in class five. The remaining coordinate
problem at that stage is a closed expression for the
local frame difference attached to an arbitrary
global class-three torsor coordinate; the obstruction
test itself is exact on actual cocycles. Higher stages
and the unrestricted global problem remain separate.
Replacing the entire module by successive
rational Tate layers computes F.2 only with its
finite-ramification and crystalline conditions; it
discards integral torsion and full finite monodromy.
For the compact pro-$\ell$ tower, F.15 supplies the missing
inverse-limit argument:
\[
 H^1_{\rm cts}(G_k,P^{(\ell)})
 \simeq\varprojlim_n H^1_{\rm cts}
       (G_k,P^{(\ell)}/\overline{\gamma_{n+1}P^{(\ell)}}).
\]
Compatible representatives exist and their compact stabilizers
cause no further obstruction. F.15 also proves the corresponding
statement for the finite-dimensional rational BCH towers with
the specified algebraic continuous action. The remaining
cohomology computations concern the higher global stages
and the full twisted Alexander module, with the chosen
arithmetic local conditions. F.21 shows concretely why
nonemptiness at every finite coefficient level need not
provide a compatible family of section classes.
For a noncompact pro-discrete
realization of $\mathcal V_{\rm ar}$, the inverse-limit comparison
requires a separate argument: F.15 gives an explicit counterexample
to injectivity without compactness or the fixed-group condition,
even with surjective transition maps.
Thus the unresolved work is the cohomology computation and,
for coefficient systems outside F.15's hypotheses, the specified
topological realization and inverse-limit comparison.
F.4 also shows why separation
of cusps cannot be inferred from the abelian calculation.
No unsolved cohomology group in this paragraph is
presented as a computed set.

\begin{statement}{F.13}{Proved}{What the normalizer version does and does not supply}
Replacing the arithmetic image by a Teichm\"uller
normalizer does not by itself define a Selmer problem.
Suppose an action on $U$ extends from an algebraic
arithmetic image $A$ to a proalgebraic normalizer $N$.
Then restriction defines
\[
 H^1_{\rm alg}(N,U)\longrightarrow H^1_{\rm alg}(A,U).
\]
For a faithfully flat image map $G\to A$ through which
the action factors, inflation is injective and its
image consists precisely of cocycles trivial on
$\ker(G\to A)$. Neither operation supplies
unramifiedness or crystallinity at a place.
\end{statement}
\begin{proof}
Restriction is substitution along $A\to N$.
For inflation, two inflated cocycles that differ by
a $U$-coboundary already differ by the same coboundary
on $A$, by faithful flat descent. Conversely a
cocycle trivial on the kernel is constant on its
cosets; descent gives the cocycle on $A$. This is
the nonabelian inflation argument with every map
specified. A local condition refers to a decomposition
group, inertia, or the $p$-adic comparison functors.
An abstract normalizer and its action contain none
of these choices. Consequently the normalizer
calculation in B gives symmetry and compatibility
constraints and possible transports such as F.5;
it cannot replace $G_k$ in the localization diagram
of F.3 without further arithmetic data.
\end{proof}

\begin{statement}{F.14}{Proved}{Lemmas that can be passed to the point-section problem}
The following conclusions apply to a proposed
point-section comparison only with their hypotheses.
A point or rational tangential point gives a
compatible section at every characteristic quotient.
It restricts to a chosen finite cover exactly under
F.11. If a unipotent quotient with the actual
arithmetic action satisfies the local hypotheses
of F.3 or F.8 and a strict global--local dimension
inequality, its integral-point locus is finite.
These facts do not prove that every section is
geometric, or that a finite-level Kummer map is
injective on all point and cusp packets.
\end{statement}
\begin{proof}
Functoriality of the étale fundamental group gives
the first assertion and compatibility of its
quotients. The cover assertion is F.11.
The finiteness assertion is exactly the localization
and analytic argument of F.3, with the refined
elliptic local conditions when F.8 is used. F.4
is an explicit counterexample to finite-level
cuspidal injectivity. The reverse implication
from arbitrary sections to points is not part
of any of these arguments. Thus these are
proved lemmas with the specified action,
base point, local conditions, and dimension
hypotheses; they are not an assumption of
the section conjecture in LM6 Part E.
\end{proof}

\begin{statement}{F.15}{Proved}{Continuous section classes and inverse limits}
Let $G$ be a topological group, let $(P_n,q_n)_{n\ge1}$ be an
inverse sequence of topological groups with continuous $G$-actions
and continuous equivariant surjections
$q_n:P_{n+1}\twoheadrightarrow P_n$, and give
$P=\varprojlim_nP_n$ its inverse-limit topology.
The natural map of pointed sets
\[
 H^1_{\rm cts}(G,P)\longrightarrow
       \varprojlim_n H^1_{\rm cts}(G,P_n)                 \tag{F.15}
\]
is bijective in either of the following cases:
\begin{enumerate}
\item Every $P_n$ is compact Hausdorff. In particular this applies
to a pro-$\ell$ group and its closed lower-central quotients.
The nonabelian $\varprojlim{}^1$ of the associated twisted
$H^0$ groups is a singleton, without requiring surjectivity
on those fixed groups.
\item For a field $k=\mathbb Q_\ell$, each
$P_n=\exp(\mathfrak p_n)$ is the BCH group of a finite-dimensional
nilpotent $k$-Lie algebra. The transition maps come from
surjective $k$-Lie homomorphisms, and $G$ acts continuously
by $k$-linear Lie automorphisms. This includes the rational
unipotent Malcev towers used here. In this case twisted
fixed groups have the Mittag--Leffler property, although
the coefficient groups need not be compact.
\end{enumerate}
Surjectivity of (F.15) in fact holds without either additional
condition. Injectivity can fail for noncompact pro-discrete
coefficient groups with surjective transition maps.
\end{statement}
\begin{proof}
We use the convention of F.1:
$c^u(g)=u^{-1}c(g)g(u)$ for the conjugate of a cocycle.
Given a compatible sequence of classes, choose a continuous
representative $c_1$. Suppose $c_n$ has been chosen. A representative
$b_{n+1}$ of the next class satisfies
$q_n b_{n+1}=c_n^{a_n}$ for some $a_n\in P_n$.
Choose a lift $\widetilde a_n\in P_{n+1}$ and put
$c_{n+1}=b_{n+1}^{\widetilde a_n^{-1}}$. Then
$q_nc_{n+1}=c_n$ exactly. Constant conjugation preserves
continuity. The map $c(g)=(c_n(g))_n$ is continuous because each
coordinate is continuous and $P$ has the inverse-limit topology;
its cocycle identity holds coordinatewise. This proves surjectivity
using only the stated surjections, with no continuous choice of
lifts of varying group elements required.

For injectivity, let $c,d:G\to P$ be continuous cocycles whose
classes agree at every level. Put
\[
 X_n=\{u\in P_n:d_n(g)=u^{-1}c_n(g)g(u)
                         \text{ for every }g\in G\}.
\]
These sets are nonempty, and $q_n(X_{n+1})\subset X_n$.
In the compact case each $X_n$ is closed in $P_n$: for each fixed
$g$, its defining equality is a closed condition in a Hausdorff
group. Thus $\prod_nX_n$ is compact. The conditions
$q_n(u_{n+1})=u_n$ are closed and have the finite-intersection
property. Indeed any finite set of them is satisfied by choosing
an element at a sufficiently high level and projecting it
downwards, with arbitrary choices in the unused coordinates.
Their common intersection is nonempty. The resulting
$u=(u_n)_n\in P$ gives $d=c^u$ and proves injectivity.

This argument also proves the stated $\varprojlim{}^1$ assertion
directly. For any inverse sequence of compact Hausdorff groups
$(H_n,f_n)$ and any $(a_n)\in\prod_nH_n$, the equations
\[
                    a_n=h_n f_n(h_{n+1})^{-1}
\]
have a solution: every finite initial set can be solved backwards,
and compactness gives a simultaneous solution. These equations
say precisely that every sequence represents the neutral class
in nonabelian $\varprojlim{}^1$. Here the relevant groups are
$H_n=P_n^{G,c_n}$ for the twisted action
$g*_cu=c_n(g)g(u)c_n(g)^{-1}$. They are closed, hence compact;
their transition maps need not be surjective.

A pro-$\ell$ group is the inverse limit of its closed
lower-central quotients. The intersection of those closed terms
is trivial because every finite $\ell$-group quotient is
nilpotent. Surjectivity onto their inverse limit follows from
compactness of the nested compatible cosets. Their characteristic
definition ensures equivariance for any continuous arithmetic
action. This justifies the pro-$\ell$ specialization in F.12;
it does not identify a full profinite group with its nilpotent
quotient tower.

For the BCH case, retain the nonempty sets $X_n$. The twisted
action differentiates to the $k$-linear Lie automorphisms
\[
       v\longmapsto\operatorname{Ad}(c_n(g))\,g(v)
                          \quad(v\in\mathfrak p_n).
\]
Exponential is a functorial bijection in the BCH construction.
Consequently
$H_n=P_n^{G,c_n}=\exp(\mathfrak h_n)$, where
$\mathfrak h_n$ is the common fixed linear subspace of these
automorphisms. For fixed $n$, the images of
$\mathfrak h_m\to\mathfrak h_n$ form a descending sequence
of finite-dimensional linear subspaces and therefore stabilize.
The point-group images $H_m\to H_n$ are exactly their
exponentials, so they stabilize as well; no compactness of
$k$-points is being assumed.

The set $X_n$ is a left torsor under $H_n$. Its images from
higher levels are nested nonempty cosets of the corresponding
images of $H_m$. Once those subgroups stabilize, the nested
cosets stabilize too. Thus $(X_n)$ is a Mittag--Leffler sequence
of nonempty sets. For completeness, put $Y_n$ equal to its
eventual image at level $n$. The maps
$Y_{n+1}\to Y_n$ are surjective: choose a level high enough
that both images have stabilized, and lift any prescribed
element of $Y_n$ from that level. Recursive choices in these
nonempty surjective sets give an element of $\varprojlim X_n$.
It again conjugates $c$ to $d$. This proves injectivity and
explains precisely why finite dimensionality replaces compactness
for the rational towers.

Finally the noncompact distinction is real. Take the profinite
group $G=(\mathbb Z_2,+)$, the discrete abelian groups
\[
 P_n=\mathbb Z\oplus\mathbb Z/2^n,
 \qquad t\cdot(a,b)=(a,b+ta\bmod 2^n),
\]
and transition maps given by reduction of the second coordinate.
These actions are continuous and the transitions are surjective.
A continuous cocycle has zero first coordinate: it would otherwise
give a nonzero continuous homomorphism from the compact group
$\mathbb Z_2$ to the torsion-free discrete group $\mathbb Z$.
Its second coordinate is therefore a continuous homomorphism
$\mathbb Z_2\to\mathbb Z/2^n$. Every such homomorphism is a
coboundary, since an integer $a$ can represent any residue
modulo $2^n$. Thus $H^1_{\rm cts}(G,P_n)=0$ for all $n$.
But $P=\mathbb Z_{\rm disc}\oplus\mathbb Z_2$, and the same
calculation gives
\[
               H^1_{\rm cts}(G,P)=\mathbb Z_2/\mathbb Z.
\]
For example $t\mapsto(0,t/3)$ represents a nonzero class
trivial at every finite level. The fixed groups
$P_n^G=2^n\mathbb Z\oplus\mathbb Z/2^n$ have nonstabilizing
images from higher levels. The program \texttt{g\_f\_limit.py}
checks the coherent residues of $1/3$ and their finite
coboundary formulas through $n=12$, and all cocycle pairs
through $n=8$. It does not claim that the integer conjugators
are compatible; their failure to be compatible is precisely
the counterexample proved above. This counterexample does not assert
failure for the particular arithmetic realization
$\mathcal V_{\rm ar}$; it shows why its topology and fixed-group
condition must be checked separately.
\end{proof}
